% Generated by the blueprint tool (blueprint-pdf/1) from a blueprint snapshot; edit the
% Markdown sources, not this file.  Compile with `latexmk -pdf` or by running
% `pdflatex` twice (the second run settles the contents and the page references).
% Images, if any, are read from assets/ next to this file.
\documentclass[11pt]{article}
\usepackage[T1]{fontenc}
\usepackage[utf8]{inputenc}
\usepackage{lmodern}
\usepackage{textcomp}
\usepackage[a4paper,margin=2.5cm,headheight=14pt]{geometry}
\usepackage{amsmath,amssymb,amsthm}
\usepackage{amscd}  % KaTeX's CD diagrams; part of amsmath, so always there
\usepackage{mathtools}
\usepackage{bm}
\usepackage{graphicx}
\usepackage{xcolor}
\usepackage{array,booktabs,tabularx}
\usepackage{fancyvrb}
\usepackage[normalem]{ulem}
\usepackage{microtype}
% Optional, and only with their fonts: a small TeX installation has the
% packages without the fonts, which fails at the first formula.
\IfFileExists{mathrsfs.sty}{\IfFileExists{rsfs10.tfm}{\usepackage{mathrsfs}}{}}{}
\IfFileExists{stmaryrd.sty}{\IfFileExists{stmary10.tfm}{\usepackage{stmaryrd}}{}}{}
\IfFileExists{cancel.sty}{\usepackage{cancel}}{}
\usepackage[hyphens]{url}
\usepackage[pdfusetitle,bookmarksnumbered=false]{hyperref}
\usepackage{bookmark}
\definecolor{bplink}{HTML}{1F4E9C}
\definecolor{bpmuted}{HTML}{5F676E}
\hypersetup{colorlinks=true,linkcolor=bplink,urlcolor=bplink,citecolor=bplink}
\setlength{\parindent}{0pt}
\setlength{\parskip}{0.55em plus 0.2em minus 0.1em}
% Lists as tight as the page's.
\makeatletter
\def\@listi{\leftmargin\leftmargini\parsep 0pt\topsep 0.3em plus 0.1em\itemsep 0.1em plus 0.1em}
\let\@listI\@listi
\def\@listii{\leftmargin\leftmarginii\labelwidth\leftmarginii\advance\labelwidth-\labelsep\parsep 0pt\topsep 0.1em\itemsep 0.1em}
% An ordered list that does not start at 1.
\newcommand{\bpstartat}[1]{\setcounter{\@enumctr}{\numexpr#1-1\relax}}
\makeatother
\setcounter{tocdepth}{3}
\emergencystretch=3em
\allowdisplaybreaks

% A numbered statement, set as a paper sets one: "Theorem 2.1 (Title)."
% #1 is the lead ("Theorem~2.1"), #2 what follows it in roman.
\newenvironment{bpstatement}[2]{\par\addvspace{0.6\baselineskip}\noindent\textbf{#1}#2\textbf{.}\enspace\ignorespaces}
  {\par\addvspace{0.4\baselineskip}}
\newenvironment{bpstatementit}[2]{\begin{bpstatement}{#1}{#2}\itshape\ignorespaces}{\end{bpstatement}}
% A step: a relation with prose of its own, read under the object it starts from.
\newenvironment{bpstep}[1]{\par\addvspace{0.3\baselineskip}\begin{list}{}{\leftmargin=1.6em\rightmargin=0pt\topsep=0pt}\item\relax
  {\small\color{bpmuted}#1}\par\nopagebreak}{\end{list}\par}
% Notes about where an object is written out.
\newenvironment{bpnote}{\par\small\color{bpmuted}\itshape}{\par}
\newcommand{\bpomitted}[1]{\par{\small\color{bpmuted}\itshape #1\par}}
\newcommand{\bpbodyheading}[1]{\par\addvspace{0.4\baselineskip}\noindent\textbf{#1}\par\nopagebreak}
\newcommand{\bprule}{\par\medskip\noindent\hfil\rule{0.3\linewidth}{0.4pt}\hfil\par\medskip}
\newcommand{\bpboundary}[1]{\par\nopagebreak{\small\color{bpmuted}#1\par}\nopagebreak}
% A reference to something this export leaves out, and one that resolves nowhere.
\newcommand{\bpout}[1]{#1\textsuperscript{\hyperlink{bp-outside}{\ensuremath{\ast}}}}
\newcommand{\bpbroken}[1]{\textcolor{red!70!black}{[#1?]}}
\newcommand{\bpslug}[1]{\texttt{#1}}
% A bibliography entry: label, anchor, entry; the label hangs in the margin.
\newcommand{\bpbibitem}[3]{\par\setbox0\hbox{[#1]\enspace}\ifdim\wd0<5em\setbox0\hbox to 5em{[#1]\hfil}\fi
  \hangindent=\wd0 \hangafter=1 \noindent\box0\phantomsection\label{#2}#3\par\smallskip}
% A character the export knows no way to set.
\newcommand{\bpglyph}[1]{\fbox{\scriptsize #1}}
% Images, at their own size unless they are wider than the line.
\newsavebox{\bpimagebox}
\newcommand{\bpimage}[1]{\sbox{\bpimagebox}{\includegraphics{#1}}%
  \ifdim\wd\bpimagebox>\linewidth\resizebox{\linewidth}{!}{\usebox{\bpimagebox}}\else\usebox{\bpimagebox}\fi}
\newcommand{\bpmissingimage}[1]{\fbox{\parbox{0.8\linewidth}{\centering\small\color{bpmuted}Image not included: #1}}}
\DefineVerbatimEnvironment{bpverbatim}{Verbatim}{fontsize=\small,xleftmargin=1em}
\renewcommand{\qedsymbol}{\ensuremath{\blacksquare}}
\hypersetup{pdftitle={Weil Conjectures — 10 Verdier and Poincaré duality},pdfsubject={Chapter 10 Verdier and Poincaré duality; all levels},pdfcreator={blueprint}}
% KaTeX commands the website knows
\providecommand{\R}{\mathbb{R}}\providecommand{\N}{\mathbb{N}}\providecommand{\Z}{\mathbb{Z}}
\providecommand{\Q}{\mathbb{Q}}\providecommand{\C}{\mathbb{C}}
\providecommand{\Reals}{\mathbb{R}}\providecommand{\reals}{\mathbb{R}}\providecommand{\Complex}{\mathbb{C}}
\providecommand{\cnums}{\mathbb{C}}\providecommand{\natnums}{\mathbb{N}}\providecommand{\Bbb}[1]{\mathbb{#1}}
\providecommand{\bold}[1]{\mathbf{#1}}\providecommand{\frak}[1]{\mathfrak{#1}}\providecommand{\mathbbm}[1]{\mathbb{#1}}
\providecommand{\lang}{\langle}\providecommand{\rang}{\rangle}\providecommand{\bra}[1]{\langle #1\rvert}
\providecommand{\ket}[1]{\lvert #1\rangle}\providecommand{\braket}[1]{\langle #1\rangle}
\providecommand{\set}[1]{\{#1\}}\providecommand{\Set}[1]{\left\{#1\right\}}
\providecommand{\infin}{\infty}\providecommand{\isin}{\in}\providecommand{\plusmn}{\pm}\providecommand{\sub}{\subset}
\providecommand{\sube}{\subseteq}\providecommand{\supe}{\supseteq}\providecommand{\exist}{\exists}
\providecommand{\bull}{\bullet}\providecommand{\larr}{\leftarrow}\providecommand{\rarr}{\rightarrow}
\providecommand{\lrarr}{\leftrightarrow}\providecommand{\Larr}{\Leftarrow}\providecommand{\Rarr}{\Rightarrow}
\providecommand{\Lrarr}{\Leftrightarrow}\providecommand{\uarr}{\uparrow}\providecommand{\darr}{\downarrow}
\providecommand{\harr}{\leftrightarrow}\providecommand{\lArr}{\Leftarrow}\providecommand{\rArr}{\Rightarrow}
\providecommand{\hArr}{\Leftrightarrow}\providecommand{\alef}{\aleph}\providecommand{\alefsym}{\aleph}
\providecommand{\real}{\Re}\providecommand{\image}{\Im}\providecommand{\weierp}{\wp}
\providecommand{\Alpha}{\mathrm{A}}\providecommand{\Beta}{\mathrm{B}}\providecommand{\Epsilon}{\mathrm{E}}
\providecommand{\Zeta}{\mathrm{Z}}\providecommand{\Eta}{\mathrm{H}}\providecommand{\Iota}{\mathrm{I}}
\providecommand{\Kappa}{\mathrm{K}}\providecommand{\Mu}{\mathrm{M}}\providecommand{\Nu}{\mathrm{N}}
\providecommand{\Omicron}{\mathrm{O}}\providecommand{\Rho}{\mathrm{P}}\providecommand{\Tau}{\mathrm{T}}
\providecommand{\Chi}{\mathrm{X}}\providecommand{\omicron}{o}\providecommand{\gt}{>}\providecommand{\lt}{<}
\providecommand{\And}{\&}\providecommand{\hdashline}{\hline}\providecommand{\operatornamewithlimits}[1]{\operatorname*{#1}}
\providecommand{\llbracket}{[\![}\providecommand{\rrbracket}{]\!]}\providecommand{\cancel}[1]{#1}
\providecommand{\mathscr}{\mathcal}
% the characters beyond ASCII this document uses
\DeclareUnicodeCharacter{2160}{\textup{I}}
\DeclareUnicodeCharacter{2161}{\textup{II}}
\DeclareUnicodeCharacter{2162}{\textup{III}}
% the project's macros ([katex.macros] in blueprint.toml)
\providecommand{\A}{}\renewcommand{\A}{\mathbf A}
\providecommand{\AEn}{}\renewcommand{\AEn}{A_{E,n}}
\providecommand{\Aaf}{}\renewcommand{\Aaf}{\mathbf A}
\providecommand{\Ab}{}\renewcommand{\Ab}{\mathsf{Ab}}
\providecommand{\AffSch}{}\renewcommand{\AffSch}{\operatorname{AffSch}}
\providecommand{\Aone}{}\renewcommand{\Aone}{\mathbf A^1}
\providecommand{\Bl}{}\renewcommand{\Bl}{\operatorname{Bl}}
\providecommand{\Br}{}\renewcommand{\Br}{\mathrm{Br}}
\providecommand{\C}{}\renewcommand{\C}{\mathbb C}
\providecommand{\Cech}{}\renewcommand{\Cech}{\check{H}}
\providecommand{\Coext}{}\renewcommand{\Coext}{\operatorname{Coext}}
\providecommand{\Cone}{}\renewcommand{\Cone}{\operatorname{Cone}}
\providecommand{\Cons}{}\renewcommand{\Cons}{\operatorname{Cons}}
\providecommand{\Cref}{}\renewcommand{\Cref}[1]{\text{#1}}
\providecommand{\D}{}\renewcommand{\D}{D}
\providecommand{\DLAFT}{}\renewcommand{\DLAFT}{\operatorname{FT}}
\providecommand{\DLAFr}{}\renewcommand{\DLAFr}{\operatorname{Fr}}
\providecommand{\DLAQ}{}\renewcommand{\DLAQ}{{\overline{\mathbf Q}_\ell}}
\providecommand{\DLRQ}{}\renewcommand{\DLRQ}{{\overline{\mathbf Q}_\ell}}
\providecommand{\Dbc}{}\renewcommand{\Dbc}{D^b_c}
\providecommand{\Dc}{}\renewcommand{\Dc}{D_{\mathrm{cons}}}
\providecommand{\Dctf}{}\renewcommand{\Dctf}{D^b_{ctf}}
\providecommand{\Dic}{}\renewcommand{\Dic}{D_{\IndCons}}
\providecommand{\Dplus}{}\renewcommand{\Dplus}{D^{+}}
\providecommand{\End}{}\renewcommand{\End}{\operatorname{End}}
\providecommand{\Et}{}\renewcommand{\Et}{\mathrm{\acute{e}t}}
\providecommand{\Ex}{}\renewcommand{\Ex}{\mathrm{Ex}}
\providecommand{\Ext}{}\renewcommand{\Ext}{\operatorname{Ext}}
\providecommand{\F}{}\renewcommand{\F}{\mathbb F}
\providecommand{\FT}{}\renewcommand{\FT}{\operatorname{FT}}
\providecommand{\Fbar}{}\renewcommand{\Fbar}{{\overline{\mathbf F}_q}}
\providecommand{\Fell}{}\renewcommand{\Fell}{\mathbf F_\ell}
\providecommand{\Fib}{}\renewcommand{\Fib}{\operatorname{Fib}}
\providecommand{\Fix}{}\renewcommand{\Fix}{\operatorname{Fix}}
\providecommand{\Fp}{}\renewcommand{\Fp}{{\mathbf F_{p}}}
\providecommand{\Fq}{}\renewcommand{\Fq}{{\mathbf F_q}}
\providecommand{\Fqn}{}\renewcommand{\Fqn}{{\mathbf F_{q^{n}}}}
\providecommand{\Frob}{}\renewcommand{\Frob}{\operatorname{Frob}}
\providecommand{\GL}{}\renewcommand{\GL}{\operatorname{GL}}
\providecommand{\Gal}{}\renewcommand{\Gal}{\operatorname{Gal}}
\providecommand{\Hens}{}\renewcommand{\Hens}{\operatorname{Hens}}
\providecommand{\Hilb}{}\renewcommand{\Hilb}{\operatorname{Hilb}}
\providecommand{\Hom}{}\renewcommand{\Hom}{\operatorname{Hom}}
\providecommand{\IndCons}{}\renewcommand{\IndCons}{\operatorname{IndCons}}
\providecommand{\Kar}{}\renewcommand{\Kar}{\operatorname{Kar}}
\providecommand{\LT}{}\renewcommand{\LT}{\operatorname{LT}}
\providecommand{\Lie}{}\renewcommand{\Lie}{\operatorname{Lie}}
\providecommand{\Loc}{}\renewcommand{\Loc}{\operatorname{Loc}}
\providecommand{\Mod}{}\renewcommand{\Mod}{\mathsf{Mod}}
\providecommand{\N}{}\renewcommand{\N}{\mathbb{N}}
\providecommand{\Needspace}{}\renewcommand{\Needspace}{\par}
\providecommand{\O}{}\renewcommand{\O}{{\mathcal{O}}}
\providecommand{\Ob}{}\renewcommand{\Ob}{\mathrm{Ob}}
\providecommand{\Oe}{}\renewcommand{\Oe}{\mathcal O_E}
\providecommand{\PP}{}\renewcommand{\PP}{\mathbf{P}}
\providecommand{\Perf}{}\renewcommand{\Perf}{\operatorname{Perf}}
\providecommand{\Pic}{}\renewcommand{\Pic}{\mathrm{Pic}}
\providecommand{\Pm}{}\renewcommand{\Pm}{\mathbf P^{m}}
\providecommand{\Pone}{}\renewcommand{\Pone}{\mathbf P^{1}}
\providecommand{\Pqn}{}\renewcommand{\Pqn}{\mathbf P^{n}}
\providecommand{\ProFin}{}\renewcommand{\ProFin}{\operatorname{ProFin}}
\providecommand{\Proet}{}\renewcommand{\Proet}{\mathrm{pro\acute{e}t}}
\providecommand{\Q}{}\renewcommand{\Q}{\mathbb Q}
\providecommand{\Qb}{}\renewcommand{\Qb}{\overline{\mathbf Q}_{\ell}}
\providecommand{\Qbar}{}\renewcommand{\Qbar}{\overline{\mathbf Q}}
\providecommand{\QbarX}{}\renewcommand{\QbarX}{\overline{\mathbf Q}_{\ell,X}}
\providecommand{\Qbarl}{}\renewcommand{\Qbarl}{{\overline{\mathbf Q}_{\ell}}}
\providecommand{\Qell}{}\renewcommand{\Qell}{{\mathbf{Q}_{\ell}}}
\providecommand{\Ql}{}\renewcommand{\Ql}{\mathbf Q_{\ell}}
\providecommand{\Qlb}{}\renewcommand{\Qlb}{{\overline{\mathbf Q}_\ell}}
\providecommand{\Qlbar}{}\renewcommand{\Qlbar}{\overline{\mathbf Q}_\ell}
\providecommand{\R}{}\renewcommand{\R}{\mathbb{R}}
\providecommand{\RG}{}\renewcommand{\RG}{\operatorname{\mathrm{R}\Gamma}}
\providecommand{\RGamma}{}\renewcommand{\RGamma}{R\Gamma}
\providecommand{\RGammac}{}\renewcommand{\RGammac}{R\Gamma_c}
\providecommand{\RHom}{}\renewcommand{\RHom}{R\mathcal{H}om}
\providecommand{\RHomX}{}\renewcommand{\RHomX}{R\mathcal H om_X}
\providecommand{\Rat}{}\renewcommand{\Rat}{\operatorname{IsRat}}
\providecommand{\Res}{}\renewcommand{\Res}{\operatorname{Res}}
\providecommand{\Rlim}{}\renewcommand{\Rlim}{R\!\varprojlim}
\providecommand{\SIm}{}\renewcommand{\SIm}{\operatorname{SIm}}
\providecommand{\Sch}{}\renewcommand{\Sch}{\mathsf{Sch}}
\providecommand{\Sh}{}\renewcommand{\Sh}{\operatorname{Sh}}
\providecommand{\Span}{}\renewcommand{\Span}{\operatorname{span}}
\providecommand{\Spec}{}\renewcommand{\Spec}{\operatorname{Spec}}
\providecommand{\Tr}{}\renewcommand{\Tr}{\operatorname{Tr}}
\providecommand{\Trsm}{}\renewcommand{\Trsm}{\operatorname{Tr}^{\mathrm{sm}}}
\providecommand{\Wt}{}\renewcommand{\Wt}{\operatorname{Wt}}
\providecommand{\Z}{}\renewcommand{\Z}{\mathbb{Z}}
\providecommand{\Zar}{}\renewcommand{\Zar}{\mathrm{zar}}
\providecommand{\Zell}{}\renewcommand{\Zell}{{\mathbf{Z}_{\ell}}}
\providecommand{\Zl}{}\renewcommand{\Zl}{{\mathbf Z_\ell}}
\providecommand{\affet}{}\renewcommand{\affet}[1]{{#1}^{\mathrm{aff}}_{\mathrm{ét}}}
\providecommand{\affproet}{}\renewcommand{\affproet}[1]{{#1}^{\mathrm{aff}}_{\mathrm{proét}}}
\providecommand{\boxt}{}\renewcommand{\boxt}{\boxtimes_A^{\mathrm L}}
\providecommand{\cF}{}\renewcommand{\cF}{\mathcal{F}}
\providecommand{\cHom}{}\renewcommand{\cHom}{R\mathcal{H}\!om}
\providecommand{\cI}{}\renewcommand{\cI}{\mathcal{I}}
\providecommand{\cL}{}\renewcommand{\cL}{\mathcal L}
\providecommand{\cM}{}\renewcommand{\cM}{\mathcal M}
\providecommand{\cN}{}\renewcommand{\cN}{\mathcal N}
\providecommand{\cO}{}\renewcommand{\cO}{\mathcal O}
\providecommand{\cT}{}\renewcommand{\cT}{\mathcal T}
\providecommand{\checked}{}\renewcommand{\checked}{\color{red}\aftergroup\restorecolor}
\providecommand{\coeff}{}\renewcommand{\coeff}{\operatorname{coeff}}
\providecommand{\coker}{}\renewcommand{\coker}{\operatorname{coker}}
\providecommand{\colim}{}\renewcommand{\colim}{\operatorname{colim}}
\providecommand{\constructedby}{}\renewcommand{\constructedby}[1]{\par\smallskip\noindent\textcolor{purple!60!black}{\textsf{Constructed by:} }
  \footnotesize\sloppy\nolinkurl{#1}\par}
\providecommand{\deepresult}{}\renewcommand{\deepresult}[1]{\par\smallskip\noindent\textcolor{orange!45!black}{
  \textsf{Blueprint role: named deep theorem node; see the cluster proof architecture.}}
  \par\noindent\textcolor{orange!45!black}{\textsf{Mathematical source:} }#1\par}
\providecommand{\dt}{}\renewcommand{\dt}{\mathbin{\otimes^{\mathbf L}}}
\providecommand{\dw}{}\renewcommand{\dw}{\operatorname{dw}}
\providecommand{\et}{}\renewcommand{\et}[1]{{#1}_{\mathrm{ét}}}
\providecommand{\ev}{}\renewcommand{\ev}{\operatorname{ev}}
\providecommand{\fib}{}\renewcommand{\fib}{\operatorname{fib}}
\providecommand{\formalizationnote}{}\renewcommand{\formalizationnote}[1]{\begin{quote}\small\raggedright\textbf{Formalization note.} #1\end{quote}}
\providecommand{\fpqc}{}\renewcommand{\fpqc}{X_{\mathrm{fpqc}}}
\providecommand{\gammacont}{}\renewcommand{\gammacont}{\Gamma_{\mathrm{cont}}}
\providecommand{\gm}{}\renewcommand{\gm}{\mathbf G_m}
\providecommand{\hocolim}{}\renewcommand{\hocolim}{\operatorname{hocolim}}
\providecommand{\holim}{}\renewcommand{\holim}{\operatorname{holim}}
\providecommand{\hrB}{}\renewcommand{\hrB}{\@ifnextchar]{\hrEnd}{\hrStep}}
\providecommand{\hrEnd}{}\renewcommand{\hrEnd}{\@firstoftwo{}}
\providecommand{\hrStep}{}\renewcommand{\hrStep}{\@firstoftwo{\hrB}}
\providecommand{\id}{}\renewcommand{\id}{\operatorname{id}}
\providecommand{\kunFl}{}\renewcommand{\kunFl}{\mathbf F_\ell}
\providecommand{\kunZn}{}\renewcommand{\kunZn}{\mathbf Z/n\mathbf Z}
\providecommand{\kuncd}{}\renewcommand{\kuncd}{\operatorname{cd}}
\providecommand{\kundim}{}\renewcommand{\kundim}{\dim}
\providecommand{\kunotimes}{}\renewcommand{\kunotimes}[1]{\otimes^{\mathbf L}_{#1}}
\providecommand{\kunsupp}{}\renewcommand{\kunsupp}{\operatorname{Supp}}
\providecommand{\lotimes}{}\renewcommand{\lotimes}{\otimes^{\mathrm L}}
\providecommand{\m}{}\renewcommand{\m}{{\mathfrak{m}}}
\providecommand{\n}{}\renewcommand{\n}{{\mathfrak{n}}}
\providecommand{\nagataGap}{}\renewcommand{\nagataGap}[1]{\par\noindent\textbf{Formalization gap.} #1\par}
\providecommand{\openobligation}{}\renewcommand{\openobligation}[1]{\par\smallskip\noindent\textcolor{red!65!black}{\textsf{Open obligation:} }
  \footnotesize\sloppy\nolinkurl{#1}\par}
\providecommand{\p}{}\renewcommand{\p}{{\mathfrak{p}}}
\providecommand{\pr}{}\renewcommand{\pr}{\operatorname{pr}}
\providecommand{\preshv}{}\renewcommand{\preshv}[1]{\mathrm{PreShv}(#1)}
\providecommand{\pro}{}\renewcommand{\pro}[1]{\mathrm{Pro}(#1)}
\providecommand{\proet}{}\renewcommand{\proet}{X_{\mathrm{proét}}}
\providecommand{\q}{}\renewcommand{\q}{{\mathfrak{q}}}
\providecommand{\rad}{}\renewcommand{\rad}{\operatorname{rad}}
\providecommand{\rank}{}\renewcommand{\rank}{\operatorname{rank}}
\providecommand{\red}{}\renewcommand{\red}{\color{red}}
\providecommand{\restorecolor}{}\renewcommand{\restorecolor}{\color{black}}
\providecommand{\retiredobligation}{}\renewcommand{\retiredobligation}[1]{\par\smallskip\noindent\textcolor{red!75!black}{\textsf{Retired invalid route:} }
  \footnotesize\sloppy\nolinkurl{#1}\par}
\providecommand{\rk}{}\renewcommand{\rk}{\operatorname{rk}}
\providecommand{\sep}{}\renewcommand{\sep}{\textnormal{sep}}
\providecommand{\shv}{}\renewcommand{\shv}[1]{\mathrm{Shv}(#1)}
\providecommand{\spec}{}\renewcommand{\spec}[1]{\mathrm{Spec}(#1)}
\providecommand{\stacks}{}\renewcommand{\stacks}[1]{\href{https://stacks.math.columbia.edu/tag/#1}{\textnormal{Stacks, Tag~#1}}}
\providecommand{\stacksproject}{}\renewcommand{\stacksproject}[1]{{\cite[\href{https://stacks.math.columbia.edu/tag/#1}{Tag #1}]{stacks-project}}}
\providecommand{\statusnote}{}\renewcommand{\statusnote}[1]{\begin{quote}\small\raggedright\textbf{Status.} #1\end{quote}}
\providecommand{\tildeideal}{}\renewcommand{\tildeideal}[2]{\widetilde{(#1, #2)}}
\providecommand{\tildering}{}\renewcommand{\tildering}[3]{{#1}_{\widetilde{(#2, #3)}}}
\providecommand{\tr}{}\renewcommand{\tr}{\operatorname{tr}}
\providecommand{\wt}{}\renewcommand{\wt}{\operatorname{wt}}

\begin{document}
\pagestyle{plain}
\begin{center}
{\LARGE\bfseries Weil Conjectures\par}
\vspace{0.7em}
{\large Chapter 10: Verdier and Poincaré duality\par}
\vspace{0.5em}
{\small\color{bpmuted}Snapshot 2026-09-26, revision 8e4cb88e2a\par}
\end{center}
\vspace{0.5em}
{\small\tableofcontents}
\bigskip\hrule\bigskip

% object: chap-Verdier-and-Poincare-duality (blueprint/chap-Verdier-and-Poincare-duality/_section.md)
\section*{\texorpdfstring{10\quad Verdier and Poincaré duality}{10 Verdier and Poincaré duality}}
\phantomsection\label{bp1}%
\addcontentsline{toc}{section}{\texorpdfstring{10\quad Verdier and Poincaré duality}{10 Verdier and Poincaré duality}}
% text: prose-10-01
\phantomsection\label{bp2}%
The main estimate in the proof of the Deligne–Laumon theorem, the goal of Part III, is an upper bound: every Frobenius eigenvalue on \(H^i\) of a smooth proper variety has absolute value at most \(q^{i/2}\). The lower bound comes from Weil's axiom (c). Suppose that the cup product pairs \(H^i\) with \(H^{2d-i}\) perfectly and that Frobenius multiplies the pairing by \(q^d\). Then the eigenvalues \(\alpha\) in degree \(i\) correspond to the eigenvalues \(q^d/\alpha\) in degree \(2d-i\), so the upper bound for \(q^d/\alpha\) is a lower bound for \(\alpha\) (Chapter \bpout{19}). The same relation gives the functional equation of the zeta function (Corollary \bpout{12.33}). This chapter proves axiom (c).

For the elliptic curve \(y^2=x^3-x\) over \(\mathbf F_3\) of the introduction, duality says that the two eigenvalues on \(H^1\) have product \(3\) (Example \hyperref[bp69]{10.23}), but it does not bound them separately.

For a compact oriented manifold \(M\) of dimension \(n\) the model is Verdier's. Let \(a:M\to\mathrm{pt}\) be the map to a point. Then \(\omega_M=a^!\mathbf Z\) is the orientation sheaf placed in degree \(-n\), and the functor \(\mathbb D=\cHom(-,\omega_M)\) is an involution on constructible complexes. Duality is the isomorphism \(R\Gamma(M,\mathbb DK)\simeq R\!\operatorname{Hom}(R\Gamma_c(M,K),\mathbf Z)\).

We follow the same route. For a scheme \(X\) with structure map \(a_X:X\to s=\operatorname{Spec}k\), put \(\omega_X=a_X^!\Lambda_s\) and \(\mathbb D_X=\cHom(-,\omega_X)\) (Definition \hyperref[bp8]{10.2}). The adjunction \(Ra_{X!}\dashv a_X^!\) gives for free a pairing of \(R\Gamma_c(X,K)\) with \(R\Gamma(X,\mathbb D_XK)\), and it also gives two exchange formulas

\[
\mathbb D_YRf_!\simeq Rf_*\mathbb D_X,\qquad
f^!\mathbb D_Y\simeq\mathbb D_Xf^*
\]

(Lemma \hyperref[bp10]{10.3}). Over an algebraically closed field the first formula, for \(a_X\), is global duality, \(R\Gamma(X,\mathbb D_XK)\simeq R\!\operatorname{Hom}_\Lambda(R\Gamma_c(X,K),\Lambda)\) (Theorem \hyperref[bp32]{10.10}). For smooth \(X\) of pure dimension \(d\), smooth purity gives \(\omega_X\simeq\Lambda_X(d)[2d]\), and global duality becomes Poincaré duality: the pairing

\[
R\Gamma_c(X,K)\otimes^{\mathbf L}R\Gamma(X,K^\vee(d))[2d]
\longrightarrow\Lambda,\qquad K^\vee=\cHom(K,\Lambda_X),
\]

built from evaluation and the trace map, is perfect (Theorem \hyperref[bp37]{10.12}).

For \(X\) smooth, proper and connected and \(\Lambda\) a field, this will give a perfect cup product pairing \(H^i\times H^{2d-i}\to\Lambda(-d)\) compatible with Frobenius (Corollary \hyperref[bp52]{10.17} and Theorem \hyperref[bp65]{10.21} in Section \hyperref[bp54]{10.5}).

What the adjunction does not give is that the pairing is perfect. For that we need Verdier biduality: the map \(K\to\mathbb D_X\mathbb D_XK\) is an isomorphism for every \(K\in D^b_{\mathrm{ctf}}(X,\Lambda)\) and every admissible \(\Lambda\) (Theorem \hyperref[bp25]{10.8}). Then \(\mathbb D_X\) is an anti-equivalence, and applying the first two exchange formulas to \(\mathbb D_XK\) gives the other two, \(\mathbb D_YRf_*\simeq Rf_!\mathbb D_X\) and \(\mathbb D_Xf^!\simeq f^*\mathbb D_Y\) (Corollary \hyperref[bp28]{10.9}).

Biduality cannot be proved by dévissage, because dualizing a piece \(t_!P\), with \(P\) locally constant on a smooth stratum \(t:T\hookrightarrow X\), twice needs the exchange formula \(\mathbb D_XRt_*\simeq t_!\mathbb D_T\), which is itself a consequence of biduality (Section \hyperref[bp12]{10.2}). We argue instead by induction on dimension, restricting to the geometric generic fibre of a function, where duality commutes with restriction up to the twist \((1)[2]\) (Lemmas \hyperref[bp19]{10.6} and \hyperref[bp22]{Lemma 10.7}).

Global duality was announced by Verdier and proved by Deligne [\hyperref[bp70]{SGA4}, XVIII]. Biduality for constructible complexes over a field is due to Deligne [\hyperref[bp71]{SGA4½}, Th. finitude].

A perfect pairing of complexes induces perfect pairings of cohomology modules over a field (Corollary \hyperref[bp40]{10.13}), but not over \(\mathcal O_E\) (Remark \hyperref[bp43]{10.14}). For this reason we formulate the duality statements in the derived category, as perfect pairings of perfect complexes (Definition \hyperref[bp35]{10.11}).

% object: sec-Exchange-formulas (blueprint/chap-Verdier-and-Poincare-duality/sec-Exchange-formulas/_section.md), content hidden
\subsection*{\texorpdfstring{10.1\quad Exchange formulas}{10.1 Exchange formulas}}
\phantomsection\label{bp3}%
\addcontentsline{toc}{subsection}{\texorpdfstring{10.1\quad Exchange formulas}{10.1 Exchange formulas}}
% text: prose-10-02
\phantomsection\label{bp4}%
After defining the dualizing complex \(\omega_X\), the Verdier dual \(\mathbb D_X\) and the double-dual map \(\delta_K:K\to\mathbb D_X\mathbb D_XK\), we prove the two exchange formulas that need only adjunctions (Lemma \hyperref[bp10]{10.3}). They say that \(\mathbb D\) turns \(Rf_!\) into \(Rf_*\) and \(f^*\) into \(f^!\). For example, for \(f=a_X\) and \(K=\Lambda_X\) the first says that \(R\Gamma(X,\omega_X)\) is the dual of \(R\Gamma_c(X,\Lambda)\), as for the homology of a manifold. The other half, with the roles of \(!\) and \(*\) reversed, needs biduality.

The induction in the proof of Theorem \hyperref[bp25]{10.8} ends on the geometric point, the case \(\dim X=0\). There we want duality to be the duality of finite projective modules, and for this constructible complexes must be perfect. Over \(\mathbf Z/\ell^n\) this is a real condition; for instance, the module \(\mathbf Z/\ell\) has infinite projective dimension over \(\mathbf Z/\ell^2\). This is why we require the coefficients to have finite Tor dimension. With this requirement, Lemma \hyperref[bp5]{10.1} shows that \(D^b_{\mathrm{ctf}}\) of the point is exactly the category of perfect complexes.

% object: lemma-perfect-complexes-geometric-point (blueprint/chap-Verdier-and-Poincare-duality/sec-Exchange-formulas/lemma-perfect-complexes-geometric-point/_section.md)
\phantomsection\label{bp5}%
\begin{bpstatementit}{Lemma~10.1}{ (Perfect complexes at the geometric point)}
Suppose \(k\) is algebraically closed. The category \(D^b_{\mathrm{ctf}}(s,\Lambda)\) identifies with \(\operatorname{Perf}(\Lambda)\), the category of bounded complexes of finite projective \(\Lambda\)-modules up to quasi-isomorphism. For \(K\in D^b_{\mathrm{ctf}}(X,\Lambda)\), both

\[
R\Gamma_c(X,K)=Ra_{X!}K,\qquad R\Gamma(X,K)=Ra_{X*}K
\]

and every geometric stalk \(K_{\bar x}\) are perfect complexes. Their duals \(R\!\operatorname{Hom}_\Lambda(-,\Lambda)\) are perfect, and the evaluation map \(P\to R\!\operatorname{Hom}_\Lambda(R\!\operatorname{Hom}_\Lambda(P,\Lambda),
\Lambda)\) is an isomorphism for each of them.
\end{bpstatementit}

% proof: proof/lemma-perfect-complexes-geometric-point
\begin{proof}[{Proof}]\phantomsection\label{bp6}%
Once \(D^b_{\mathrm{ctf}}(s,\Lambda)\) is identified with \(\operatorname{Perf}(\Lambda)\), the other assertions follow, because the direct images and the stalks of \(K\) are objects of \(D^b_{\mathrm{ctf}}(s,\Lambda)\).

By the comparisons of Section \bpout{2.13}, taking global sections identifies \(D^b_{\mathrm{ctf}}(s,\Lambda)\) with the full subcategory of \(D(\Lambda)\) of complexes with bounded finitely generated cohomology and finite Tor dimension, and we show that these complexes are perfect separately for the three kinds of coefficients. A constructible complex with finite coefficients has bounded finite cohomology and finite Tor amplitude, so it is perfect: take a bounded-above resolution by finite free modules and truncate it below the Tor bound. Then the last syzygy that remains is finitely presented and flat, hence projective, so the truncated resolution is a bounded complex of finite projective modules. With integral coefficients, by contrast, the constructible complexes on the point are those with bounded finitely generated \(\mathcal O_E\)-cohomology, and since a discrete valuation ring has global dimension one, these complexes are perfect. Finally, with field coefficients, bounded complexes with finite-dimensional cohomology are perfect.

By Proposition \bpout{7.1}, the two direct images of \(a_X\) lie in \(D^b_{\mathrm{ctf}}(s,\Lambda)\) and are therefore perfect. If we represent such a complex by a bounded complex of finite projective modules, its termwise dual is again of this form, and the double evaluation map is termwise an isomorphism. The stalk \(K_{\bar x}\) is the complex of global sections of \(\bar x^*K\in D^b_{\mathrm{ctf}}(s,\Lambda)\) (Proposition \bpout{7.1}), so it is perfect as well.
\end{proof}

% text: prose-10-03
\phantomsection\label{bp7}%
On the point, then, duality is the linear duality of perfect complexes, and the double dual is the identity. On a general \(X\), however, we dualize into \(a_X^!\Lambda_s\) instead of \(\Lambda\). By the adjunction \(Ra_{X!}\dashv a_X^!\), maps into \(a_X^!\Lambda_s\) are the same as linear forms on compactly supported cohomology.

% object: definition-verdier-dual-evaluation (blueprint/chap-Verdier-and-Poincare-duality/sec-Exchange-formulas/definition-verdier-dual-evaluation/_section.md)
\phantomsection\label{bp8}%
\begin{bpstatement}{Definition~10.2}{ (Verdier dual and evaluation)}
The \emph{dualizing complex} and the \emph{Verdier dual} are

\[
\omega_X=a_X^!\Lambda_s,\qquad
\mathbb D_XK=\cHom_{\Lambda_X}(K,\omega_X).
\]

Thus \(\mathbb D_X:D(X_{\mathrm{pro}\acute{e}t},\Lambda)^{\mathrm{op}}\to
D(X_{\mathrm{pro}\acute{e}t},\Lambda)\) is an exact functor. The \emph{double-dual morphism}

\[
\delta_K:K\longrightarrow\mathbb D_X\mathbb D_XK
\]

is the tensor–Hom adjoint of the composite \(K\otimes^{\mathbf L}\mathbb D_XK\simeq\mathbb D_XK\otimes^{\mathbf L}K
\to\omega_X\) of the symmetry isomorphism, with its Koszul sign, and evaluation.
\end{bpstatement}

% text: prose-10-04
\phantomsection\label{bp9}%
On the geometric point we have \(\omega_s=\Lambda_s\), so \(\delta\) is the evaluation map of Lemma \hyperref[bp5]{10.1}.

The proof of biduality uses only the exchange formulas that follow from adjunctions. We first write down the comparison maps. Let \(f:X\to Y\) be a morphism, \(K\in D^b_{\mathrm{ctf}}(X,\Lambda)\) and \(L\in D^b_{\mathrm{ctf}}(Y,\Lambda)\). The map

\[
Rf_*\cHom_X(K,f^!L)\longrightarrow\cHom_Y(Rf_!K,L)
\tag{10.1}
\]

is the tensor–Hom adjoint of the composite

\[
\begin{aligned}
Rf_*\cHom_X(K,f^!L)\otimes^{\mathbf L}Rf_!K
&\simeq Rf_!\bigl(f^*Rf_*\cHom_X(K,f^!L)\otimes^{\mathbf L}K\bigr)\\
&\longrightarrow Rf_!\bigl(\cHom_X(K,f^!L)\otimes^{\mathbf L}K\bigr)
\longrightarrow Rf_!f^!L\xrightarrow{\varepsilon_f}L
\end{aligned}
\]

of the projection formula, the counit of \(f^*\dashv Rf_*\), evaluation and the counit \(\varepsilon_f\) of \(Rf_!\dashv f^!\). In particular, for \(L=\omega_Y\), using \(f^!\omega_Y=f^!a_Y^!\Lambda_s=\omega_X\), it is a map \(Rf_*\mathbb D_XK\to\mathbb D_YRf_!K\). The map \(f^!\mathbb D_YL\to\mathbb D_Xf^*L\) is the tensor–Hom adjoint of the map \(f^!\mathbb D_YL\otimes^{\mathbf L}f^*L\to f^!\omega_Y=\omega_X\) that corresponds under \(Rf_!\dashv f^!\) to

\[
Rf_!(f^!\mathbb D_YL\otimes^{\mathbf L}f^*L)
\simeq Rf_!f^!\mathbb D_YL\otimes^{\mathbf L}L
\xrightarrow{\varepsilon_f\otimes1}\mathbb D_YL\otimes^{\mathbf L}L
\longrightarrow\omega_Y,
\]

the last arrow being evaluation.

% object: lemma-internal-adjunction-two-exchange-formulas (blueprint/chap-Verdier-and-Poincare-duality/sec-Exchange-formulas/lemma-internal-adjunction-two-exchange-formulas/_section.md)
\phantomsection\label{bp10}%
\begin{bpstatementit}{Lemma~10.3}{ (Internal adjunction and two exchange formulas)}
Let \(f:X\to Y\) be a morphism, \(K\in D^b_{\mathrm{ctf}}(X,\Lambda)\) and \(L\in D^b_{\mathrm{ctf}}(Y,\Lambda)\). The map (10.1) and the two maps derived from it and defined above are isomorphisms

\begin{align*}
\cHom_Y(Rf_!K,L)&\simeq
Rf_*\cHom_X(K,f^!L),\tag{10.2}\\
e_K:\mathbb D_YRf_!K&\xrightarrow{\ \sim\ } Rf_*\mathbb D_XK,
\tag{10.3}\\
f^!\mathbb D_YL&\simeq\mathbb D_Xf^*L.
\tag{10.4}\end{align*}

If \(f\) is proper, then the composite

\[
Rf_*K\xrightarrow{Rf_*\delta_K}Rf_*\mathbb D_X\mathbb D_XK
\simeq\mathbb D_YRf_*\mathbb D_XK
\]

of \(Rf_*\delta_K\) and \(e_{\mathbb D_XK}\) equals the composite

\[
Rf_*K\xrightarrow{\delta_{Rf_*K}}\mathbb D_Y\mathbb D_YRf_*K
\simeq\mathbb D_YRf_*\mathbb D_XK
\]

of \(\delta_{Rf_*K}\) and \(\mathbb D_Y\) applied to \(e_K\).
\end{bpstatementit}

% proof: proof/lemma-internal-adjunction-two-exchange-formulas
\begin{proof}[{Proof}]\phantomsection\label{bp11}%
We check each isomorphism by Yoneda. After applying \(\Hom(N,-)\), we identify both sides with the same Hom group by the adjunctions and the projection formula. We only have to make sure that the test objects \(N\) lie in the localizing category on which \(Rf_!\dashv f^!\) is defined.

The adjunction \(Rf_!\dashv f^!\) lives on the localizing categories \(\Loc(D^b_{\mathrm{ctf}}(X,\Lambda))\) and \(\Loc(D^b_{\mathrm{ctf}}(Y,\Lambda))\), and the functors in the argument preserve them. Indeed, pullback is exact, preserves direct sums, and carries the constructible generators to constructible objects by Proposition \bpout{7.1}, so it maps \(\Loc(D^b_{\mathrm{ctf}}(Y,\Lambda))\) into \(\Loc(D^b_{\mathrm{ctf}}(X,\Lambda))\). For the same reasons, tensoring with a fixed constructible complex preserves the localizing category. In particular, for \(N\in\Loc(D^b_{\mathrm{ctf}}(Y,\Lambda))\) the objects \(f^*N\otimes^{\mathbf L}K\) and \(N\otimes^{\mathbf L}Rf_!K\) lie in \(\Loc(D^b_{\mathrm{ctf}}(X,\Lambda))\) and \(\Loc(D^b_{\mathrm{ctf}}(Y,\Lambda))\), respectively.

The second step of the computation below is the projection formula

\[
Rf_!(f^*N\otimes^{\mathbf L}K)\simeq N\otimes^{\mathbf L}Rf_!K
\qquad(N\in\Loc(D^b_{\mathrm{ctf}}(Y,\Lambda))).
\]

For finite coefficients it follows from Proposition \bpout{5.15} and [Stacks, \href{https://stacks.math.columbia.edu/tag/0GL5}{0GL5}], and for rational coefficients it is Proposition \bpout{6.23}. For integral coefficients we compactify \(f=pj\), with \(j\) an open immersion and \(p\) proper, and combine Lemma \bpout{6.26} with \(j_!(M\otimes^{\mathbf L}j^*N)
 \simeq j_!M\otimes^{\mathbf L}N\).

By Proposition \bpout{7.1}, both \(\cHom_Y(Rf_!K,L)\) and \(Rf_*\cHom_X(K,f^!L)\) belong to \(D^b_{\mathrm{ctf}}(Y,\Lambda)\), so we may test the comparison map between them by Yoneda inside the full subcategory \(\Loc(D^b_{\mathrm{ctf}}(Y,\Lambda))\). For every object \(N\) of this category we obtain, by the adjunctions and the projection formula,

\[
\begin{aligned}
\Hom_Y(N,\cHom_Y(Rf_!K,L))
&\simeq\Hom_Y(N\otimes^{\mathbf L}Rf_!K,L)\\
&\simeq\Hom_Y(Rf_!(f^*N\otimes^{\mathbf L}K),L)\\
&\simeq\Hom_X(f^*N\otimes^{\mathbf L}K,f^!L)\\
&\simeq\Hom_Y(N,Rf_*\cHom_X(K,f^!L)).
\end{aligned}
\]

The composite bijection is induced by the map (10.1). By Yoneda this map is an isomorphism, which is (10.2). Taking \(L=\omega_Y\) and using \(f^!\omega_Y=\omega_X\), we get \(e_K\). For (10.4), on the other hand, note that both \(f^!\mathbb D_YL\) and \(\mathbb D_Xf^*L\) lie in \(D^b_{\mathrm{ctf}}(X,\Lambda)\). Since tensoring with \(f^*L\) preserves \(\Loc(D^b_{\mathrm{ctf}}(X,\Lambda))\), we may test on objects \(M\) of this category:

\[
\begin{aligned}
\Hom_X(M,f^!\mathbb D_YL)
&\simeq\Hom_Y(Rf_!M,\cHom_Y(L,\omega_Y))\\
&\simeq\Hom_Y(Rf_!M\otimes^{\mathbf L}L,\omega_Y)\\
&\simeq\Hom_X(M\otimes^{\mathbf L}f^*L,f^!\omega_Y)\\
&\simeq\Hom_X(M,\mathbb D_Xf^*L).
\end{aligned}
\]

Again, the bijection is induced by the map defined before the lemma, and by Yoneda we get (10.4).

The comparison of double-dual maps is used in the proof of biduality for \(f=a_X\), where it identifies \(R\Gamma(X,\delta_K)\) with the double-dual map of the perfect complex \(R\Gamma(X,K)\). Both maps are determined by pairings, and we compare the pairings. For proper \(f\) we have \(Rf_!=Rf_*\), and the isomorphism \(e_K\) is adjoint to the pairing

\[
Rf_*K\otimes^{\mathbf L}Rf_*\mathbb D_XK
\longrightarrow Rf_*(K\otimes^{\mathbf L}\mathbb D_XK)
\longrightarrow Rf_*\omega_X\longrightarrow\omega_Y.
\]

Using \(e_K\) and \(e_{\mathbb D_XK}\), we view both \(Rf_*\delta_K\) and \(\delta_{Rf_*K}\) as maps \(Rf_*K\to\mathbb D_YRf_*\mathbb D_XK\). Up to the symmetry of the tensor product, both are adjoint to this pairing. Indeed, for \(\delta_{Rf_*K}\) this holds because under \(e_K\) evaluation on \(Y\) corresponds to the pairing. Similarly, for \(Rf_*\delta_K\) it holds because \(\delta_K\) is adjoint to evaluation on \(X\) and \(e_{\mathbb D_XK}\) is adjoint to the analogous pairing.
\end{proof}


% object: sec-Verdier-biduality (blueprint/chap-Verdier-and-Poincare-duality/sec-Verdier-biduality/_section.md), content hidden
\subsection*{\texorpdfstring{10.2\quad Verdier biduality}{10.2 Verdier biduality}}
\phantomsection\label{bp12}%
\addcontentsline{toc}{subsection}{\texorpdfstring{10.2\quad Verdier biduality}{10.2 Verdier biduality}}
% text: prose-10-05
\phantomsection\label{bp13}%
On a smooth curve \(C\) over an algebraically closed field, with \(\omega_C=\Lambda_C(1)[2]\), biduality is easy to see for the two basic kinds of sheaves. Indeed, for a locally constant sheaf \(L\) of free \(\Lambda\)-modules the dual is \(L^\vee(1)[2]\), and dualizing twice gives back \(L\). For a skyscraper \(i_*M\) at a closed point \(i:x\to C\), the exchange formula \(\mathbb D_Ci_*\simeq i_*\mathbb D_x\) for the proper map \(i\) shows that the dual is \(i_*\mathbb D_xM\), so biduality reduces to biduality on the point. In general, a constructible sheaf is an extension of \(i_*M\) by \(j_!L\), where now \(i\) is the inclusion of a finite set of closed points and \(L\) is locally constant on its complement \(j:U\hookrightarrow C\). The open piece is the difficult one: we have \(\mathbb D_C(j_!L)\simeq Rj_*\mathbb D_UL\), and to dualize \(Rj_*\) we would need the second exchange formula. That formula, however, depends on biduality. In dimension one this can still be checked by hand, using the cohomology of the henselian punctured disc, whereas in higher dimension the direct images \(Rj_*\) along the boundary of a stratum have no such description.

% object: example-biduality-punctured-line (blueprint/chap-Verdier-and-Poincare-duality/sec-Verdier-biduality/example-biduality-punctured-line/_section.md)
\phantomsection\label{bp14}%
\begin{bpstatement}{Example~10.4}{ (The dual of $j_!\Lambda$ on the line)}
Let \(k\) be algebraically closed, \(X=\mathbf A^1_k\), \(i:0\to X\) the origin and \(j:U=\mathbf G_{m,k}\to X\) its complement, and take \(\Lambda\) finite. The sheaf \(j_!\Lambda_U\) has stalk zero at \(0\), but its dual does not. By (10.3) for \(j\) and \(\omega_U=\Lambda_U(1)[2]\),

\[
\mathbb D_X(j_!\Lambda_U)\simeq Rj_*\mathbb D_U\Lambda_U
=Rj_*\Lambda_U(1)[2],
\]

and by Example \bpout{6.25} its stalk at \(0\) is \(\Lambda(1)\) in degree \(-2\) and \(\Lambda\) in degree \(-1\). The term in degree \(-1\) comes from \(R^1j_*\Lambda_U=i_*\Lambda(-1)\). The same answer comes out of the triangle \(j_!\Lambda_U\to\Lambda_X\to i_*\Lambda\to\). Dualizing it, with \(\mathbb D_X\Lambda_X=\Lambda_X(1)[2]\) and \(\mathbb D_X(i_*\Lambda)=i_*\mathbb D_0\Lambda=i_*\Lambda\), gives a triangle

\[
i_*\Lambda\longrightarrow\Lambda_X(1)[2]\longrightarrow
\mathbb D_X(j_!\Lambda_U)\longrightarrow.
\]

On stalks at \(0\) the first map, from degree \(0\) to degree \(-2\), is zero. Hence the cone has \(\Lambda(1)\) in degree \(-2\) and \(\Lambda\) in degree \(-1\).

Dualizing once more, biduality predicts \(\mathbb D_X\mathbb D_X(j_!\Lambda_U)
\simeq j_!\Lambda_U\), with zero stalk at \(0\). In other words, the second dual cancels the stalk that the first created. Here this follows from the two easy cases. The map \(\delta\) is a morphism of triangles, and it is an isomorphism on \(\Lambda_X\) and on \(i_*\Lambda\), so by the five lemma it is an isomorphism on \(j_!\Lambda_U\). The argument used, however, that \(\Lambda_U\) extends to the locally constant sheaf \(\Lambda_X\). A locally constant \(L\) on \(U\) with nontrivial monodromy around \(0\) has no such extension, and that is the case which needs the henselian disc.
\end{bpstatement}

% text: prose-10-06
\phantomsection\label{bp15}%
The proof of Theorem \hyperref[bp25]{10.8} needs two preliminary facts. First, Lemma \hyperref[bp16]{10.5} shows that \(\mathbb D_X\) preserves \(D^b_{\mathrm{ctf}}(X,\Lambda)\). Second, Lemmas \hyperref[bp19]{10.6} and \hyperref[bp22]{Lemma 10.7} show that \(\mathbb D_X\) and \(\delta\) commute with restriction to a geometric generic fibre up to the twist \((1)[2]\), and this twist cancels in the double dual.

% object: lemma-verdier-duality-preserves-constructibility (blueprint/chap-Verdier-and-Poincare-duality/sec-Verdier-biduality/lemma-verdier-duality-preserves-constructibility/_section.md)
\phantomsection\label{bp16}%
\begin{bpstatementit}{Lemma~10.5}{ (Verdier duality preserves constructibility)}
For every admissible \(\Lambda\), \(\mathbb D_X\) restricts to a functor \(D^b_{\mathrm{ctf}}(X,\Lambda)^{\mathrm{op}}\to D^b_{\mathrm{ctf}}(X,\Lambda)\).
\end{bpstatementit}

% proof: proof/lemma-verdier-duality-preserves-constructibility
\begin{proof}[{Proof}]\phantomsection\label{bp17}%
The lemma follows from the stability of \(D^b_{\mathrm{ctf}}\) under the six operations. Proposition \bpout{7.1} shows that \(a_X^!\) and internal Hom preserve \(D^b_{\mathrm{ctf}}\). As \(\Lambda_s\) lies in \(D^b_{\mathrm{ctf}}(s,\Lambda)\), both \(\omega_X=a_X^!\Lambda_s\) and \(\mathbb D_XK\) lie in \(D^b_{\mathrm{ctf}}(X,\Lambda)\) for \(K\in D^b_{\mathrm{ctf}}(X,\Lambda)\). We also give a direct argument with finite coefficients and smooth strata, because the proof of Lemma \hyperref[bp19]{10.6} uses its dévissage to generators \(t_!P\) again.

\emph{Dévissage.} The complex \(K\) is built by finitely many cones from complexes \(t_!P\), for locally closed \(t:T\hookrightarrow X\) with \(T\) smooth and \(P\) a locally constant perfect complex. To see this, choose a finite stratification by smooth strata, disregarding nilpotents, on which the restrictions \(P\) of \(K\) are locally constant perfect complexes, and use the localization triangles. The bounded flat models of Lemma \bpout{3.13} and their stupid filtrations make this precise. Since \(\mathbb D_X\) is exact and \(D^b_{\mathrm{ctf}}(X,\Lambda)\) is closed under cones, it suffices to show that \(\mathbb D_X(t_!P)\) lies in \(D^b_{\mathrm{ctf}}(X,\Lambda)\).

\emph{The generators.} Internal adjunction (10.3) for \(t\) gives

\[
\mathbb D_X(t_!P)\simeq Rt_*\mathbb D_TP.
\]

On a smooth stratum of dimension \(r\), Theorem \bpout{9.3} identifies the dual \(\mathbb D_TP\) with \(P^\vee(r)[2r]\), which is locally a perfect complex. By Proposition \bpout{7.1} its ordinary direct image lies in \(D^b_{\mathrm{ctf}}(X,\Lambda)\).
\end{proof}

% text: prose-10-07
\phantomsection\label{bp18}%
We prove biduality by induction on dimension, restricting the double-dual morphism to the geometric generic fibre of a map to \(\mathbf A^1_k\). This needs comparison maps between the operations on a scheme and on its geometric generic fibre. Let \(B/k\) be a smooth integral curve and \(\bar\eta=\operatorname{Spec}\overline{k(B)}\), and for a \(B\)-scheme \(T\) of finite type let \(g_T:T_{\bar\eta}=T\times_B\bar\eta\to T\) be the projection, with pullback \(g_T^*:D(T_{\mathrm{pro}\acute{e}t},A)\to D(T_{\bar\eta,\mathrm{pro}\acute{e}t},A)\), where \(A\) is a finite coefficient ring. First, for a morphism \(h:T\to X\) of \(B\)-schemes the base change map \(g_X^*Rh_*L\to Rh_{\bar\eta,*}g_T^*L\) is adjoint to \(g_T^*\) of the counit \(h^*Rh_*L\to L\). Second, for a closed immersion \(i:Z\to X\) of \(B\)-schemes, the pullback \(i_{\bar\eta,*}g_Z^*i^!M=g_X^*i_*i^!M\to g_X^*M\) of the counit corresponds under \(i_{\bar\eta,*}\dashv i_{\bar\eta}^!\) to a map \(g_Z^*i^!M\to i_{\bar\eta}^!g_X^*M\). Finally, for internal Hom, \(g_X^*\cHom_X(K,M)\to\cHom_{X_{\bar\eta}}(g_X^*K,g_X^*M)\) is the tensor–Hom adjoint of \(g_X^*\) of evaluation \(\cHom_X(K,M)\otimes^{\mathbf L}K\to M\). The next lemma says that these maps are isomorphisms.

% object: lemma-restriction-geometric-generic-fibre (blueprint/chap-Verdier-and-Poincare-duality/sec-Verdier-biduality/lemma-restriction-geometric-generic-fibre/_section.md)
\phantomsection\label{bp19}%
\begin{bpstatementit}{Lemma~10.6}{ (Restriction to a geometric generic fibre)}
In the setting above, let \(X\) be a \(B\)-scheme of finite type, let \(h:T\to X\) be a separated morphism of \(B\)-schemes of finite type, let \(i:Z\to X\) be a closed immersion, and let \(L\) on \(T\) and \(K,M\) on \(X\) be bounded constructible complexes, with \(K\in D^b_{\mathrm{ctf}}(X,A)\). The three maps

\[
\begin{aligned}
g_X^*Rh_*L&\longrightarrow Rh_{\bar\eta,*}g_T^*L,\\
g_Z^*i^!M&\longrightarrow i_{\bar\eta}^!g_X^*M,\\
g_X^*\cHom_X(K,M)&\longrightarrow
\cHom_{X_{\bar\eta}}(g_X^*K,g_X^*M)
\end{aligned}
\]

defined above are isomorphisms.
\end{bpstatementit}

% proof: proof/lemma-restriction-geometric-generic-fibre
\begin{proof}[{Proof}]\phantomsection\label{bp20}%
We prove the comparisons in the order stated, each using the previous ones: the direct image by continuity along a limit of étale maps, \(i^!\) by the localization triangle, and internal Hom by dévissage to generators \(t_!P\). Throughout, with finite coefficients Proposition \bpout{5.15} lets us compute \(D^b_{\mathrm{ctf}}(X,A)\) and the direct images by their classical étale counterparts.

\emph{Direct image.} The geometric generic point is a limit of étale \(B\)-schemes. Indeed, a separable closure of \(k(B)\) is the inverse limit of nonempty affine opens in the normalizations of \(B\) in its finite separable extensions, with the branch loci removed, and each stage is étale over an open of \(B\). The functor \(Rh_*\) commutes with pullback to each stage, because pullback to an étale \(B\)-scheme is restriction to an object of the small site. It remains to pass to the limit.

Every affine object étale over the limit descends to a stage. Continuity of cohomology (Corollary \bpout{2.29}) therefore shows that, degree by degree, the derived sections with finite coefficients commute with this filtered limit along affine transition maps, provided the complexes vanish below one fixed degree. They do, because \(L\) is bounded. Moreover, all the direct images lie in one fixed range of degrees. The base changes of \(h\) to the stages are of finite type of bounded dimension, and [Stacks, \href{https://stacks.math.columbia.edu/tag/0F10}{0F10}] bounds their \(\ell\)-cohomological dimensions (Definition \bpout{2.45}) independently of the stage. This proves the first comparison over the separable closure. It also holds over the algebraic closure, since a purely inseparable extension induces the same classical étale topos.

\emph{Exceptional inverse image.} Let \(j\) be the open complement of \(i\) and pull back

\[
i_*i^!M\longrightarrow M\longrightarrow Rj_*j^*M
\longrightarrow.
\]

We compare the pulled-back triangle with the triangle with closed support on the generic fibre, whose first term is \(i_{\bar\eta,*}i_{\bar\eta}^!g_X^*M\). Their second terms agree. The third terms agree too, because the first comparison identifies the pullback of \(Rj_*j^*M\) with \(Rj_{\bar\eta,*}j_{\bar\eta}^*g_X^*M\), and the map between them is the same. It follows that the first terms are determined functorially. Indeed, there are no nonzero maps, in any degree, from a complex supported on \(Z_{\bar\eta}\) to one of the form \(Rj_{\bar\eta,*}N\), so the first term of each triangle is the fibre of the second map. Now pullback commutes with direct image along a closed immersion, so the first term of the pulled-back triangle is \(i_{\bar\eta,*}g_Z^*i^!M\). By construction, the induced map \(i_{\bar\eta,*}g_Z^*i^!M\to i_{\bar\eta,*}i_{\bar\eta}^!g_X^*M\) is adjoint to the pulled-back counit, that is, it is \(i_{\bar\eta,*}\) of our comparison map, and it is an isomorphism. Since \(i_{\bar\eta,*}\) is fully faithful, the comparison map itself is an isomorphism.

\emph{Internal Hom.} The complexes \(K\) for which the comparison is an isomorphism are closed under cones and shifts, because both sides are exact functors of \(K\). The dévissage of Lemma \hyperref[bp16]{10.5} therefore reduces us to \(K=t_!P\), where \(t\) is the inclusion of a stratum and \(P\) is locally perfect on it. For these generators, internal adjunction and the dualizability of the locally perfect complex \(P\) express internal Hom through operations we have already compared:

\[
\cHom_X(t_!P,M)
\simeq Rt_*(P^\vee\otimes^{\mathbf L}t^!M).
\]

We check that \(g^*\) commutes with each operation on the right. For \(t^!\) we factor \(t\) as a closed immersion followed by an open immersion: for the open immersion \(t^!\) is restriction, which commutes with \(g^*\), and for the closed immersion we use the second comparison. Similarly, pullback commutes with \(Rt_*\) by the first comparison, and it commutes with tensor products and with \(P^\vee\). The first two comparisons apply here, because by Lemma \hyperref[bp16]{10.5} all complexes used remain bounded constructible. Hence the comparison holds for \(K=t_!P\).
\end{proof}

% text: prose-10-08
\phantomsection\label{bp21}%
For \(B=\mathbf A^1_k\) the dualizing complex changes by a twist under restriction to the geometric generic fibre, and this twist cancels in the double dual.

% object: lemma-double-dual-geometric-generic-fibre (blueprint/chap-Verdier-and-Poincare-duality/sec-Verdier-biduality/lemma-double-dual-geometric-generic-fibre/_section.md)
\phantomsection\label{bp22}%
\begin{bpstatementit}{Lemma~10.7}{ (Double duals on the geometric generic fibre)}
In the setting of Lemma \hyperref[bp19]{10.6}, let \(B=\mathbf A^1_k\) and \(K\in D^b_{\mathrm{ctf}}(X,A)\). For every affine open \(V\) of \(X\) there is an isomorphism \(g_V^*\omega_V\simeq\omega_{V_{\bar\eta}}(1)[2]\), and the isomorphism \(g_V^*\mathbb D_V\mathbb D_V(K|_V)\simeq\mathbb D_{V_{\bar\eta}}
\mathbb D_{V_{\bar\eta}}g_V^*(K|_V)\) obtained from it and from the comparison for internal Hom (applied twice) carries \(g_V^*\delta_{K|_V}\) to \(\delta_{g_V^*(K|_V)}\).
\end{bpstatementit}

% proof: proof/lemma-double-dual-geometric-generic-fibre
\begin{proof}[{Proof}]\phantomsection\label{bp23}%
Replacing \(X\) by an affine open we may assume \(X\) affine. The twist \((1)[2]\) appears because \(X\) has one more dimension than its generic fibre, and we see it by embedding both in affine spaces whose dimensions differ by one. To do this, choose a closed immersion \(i\) over \(\mathbf A^1_k\) of \(X\) into \(\mathbf A^r_{\mathbf A^1_k}\). Since \(\mathbf A^r_{\mathbf A^1_k}\) is smooth of dimension \(r+1\) over \(k\), by composition of exceptional inverse images and Theorem \bpout{9.3} we get \(\omega_X\simeq i^!A(r+1)[2r+2]\). In the same way, the geometric generic fibre is a closed subscheme of \(\mathbf A^r_{\bar\eta}\), so \(\omega_{X_{\bar\eta}}\simeq i_{\bar\eta}^!A(r)[2r]\). By the second comparison, \(g_X^*i^!A(r+1)[2r+2]\simeq i_{\bar\eta}^!A(r+1)[2r+2]\), which is \(\omega_{X_{\bar\eta}}(1)[2]\).

For the double-dual maps, note that the comparison for internal Hom is adjoint to the pullback of the evaluation \(K\otimes^{\mathbf L}\cHom_X(K,M)\to M\), so it carries evaluation to evaluation. Together with the twist it gives \(g_X^*\mathbb D_XK\simeq\mathbb D_{X_{\bar\eta}}(g_X^*K)(1)[2]\). This twist does no harm, because the double-dual map does not change if the dualizing complex is replaced by an isomorphic object \(\omega'\): an isomorphism \(\omega\simeq\omega'\) is used once in each of the two duals, and the two uses cancel in \(\delta\). We apply this to \(\omega'=\omega_{X_{\bar\eta}}(1)[2]\), using that \(\cHom(\cHom(N,\omega(1)[2]),\omega(1)[2])\simeq\cHom(\cHom(N,\omega),\omega)\) compatibly with evaluation. It follows that the isomorphism \(g_X^*\mathbb D_X\mathbb D_XK\simeq\mathbb D_{X_{\bar\eta}}
\mathbb D_{X_{\bar\eta}}g_X^*K\) obtained by applying the comparison twice carries \(g_X^*\delta_K\) to \(\delta_{g_X^*K}\).
\end{proof}

% text: prose-10-09
\phantomsection\label{bp24}%
With these lemmas we can prove biduality.

% object: theorem-verdier-biduality (blueprint/chap-Verdier-and-Poincare-duality/sec-Verdier-biduality/theorem-verdier-biduality/_section.md)
\phantomsection\label{bp25}%
\begin{bpstatementit}{Theorem~10.8}{ (Verdier biduality)}
For every admissible coefficient ring \(\Lambda\) and every scheme \(X\),

\[
\delta_K:K\xrightarrow{\sim}\mathbb D_X\mathbb D_XK
\qquad(K\in D^b_{\mathrm{ctf}}(X,\Lambda)).
\]

Thus \(\mathbb D_X\) is an exact anti-equivalence of \(D^b_{\mathrm{ctf}}(X,\Lambda)\) with itself.
\end{bpstatementit}

% proof: proof/theorem-verdier-biduality
\begin{proof}[{Proof}]\phantomsection\label{bp26}%
We first treat finite coefficients \(A\) over an algebraically closed field \(k\), and deduce the other coefficient rings and a finite base field at the end. In this main case we reduce to proper \(X\), show by induction that the cone of \(\delta_K\) has zero stalks at all nonclosed points, and compute global sections to see that its stalks at the finitely many remaining closed points vanish. This induction is on \(d=\dim X\). Since the induction step passes to a geometric generic fibre over \(\overline{k(t)}\), we run it simultaneously over all algebraically closed fields of characteristic different from \(\ell\). In dimension zero, \(X_{\mathrm{red}}\) is a finite disjoint union of copies of \(\operatorname{Spec}k\), and by nilpotent invariance we may replace \(X\) by it. On each copy, moreover, \(\omega=\Lambda\) and \(\delta_K\) is the evaluation map of a perfect complex, which is an isomorphism by Lemma \hyperref[bp5]{10.1}.

We may assume that \(X\) is proper of dimension \(d\); properness is used only in the last step, to compute global sections. This reduction rests on the fact that the double-dual morphism commutes with restriction to opens. This holds because restriction to an open \(j:U\to X\) commutes with internal Hom, satisfies \(j^*\omega_X=j^!\omega_X=\omega_U\), and identifies all the evaluation maps. Given this, choose a compactification \(j:X\hookrightarrow\overline X\), replace the target by the closure of \(X\), which has dimension \(d\), and extend \(K\) by zero. Then \(j_!K\) is in \(D^b_{\mathrm{ctf}}(\overline X,A)\), and the restriction of its double-dual morphism to \(X\) is \(\delta_K\).

Let \(Q=\Cone(\delta_K)\). We have to show that \(Q=0\). Since \(Q\) is bounded and constructible by Lemma \hyperref[bp16]{10.5}, it is enough to show that all its geometric stalks vanish.

At nonclosed points the induction hypothesis applies. Let \(x\) be a nonclosed point of \(X\). We choose an affine neighbourhood \(U\) of \(x\) and an algebra generator \(t\) whose image in \(\kappa(x)\) is transcendental over \(k\). Such a generator exists because a nonclosed point has positive transcendence degree. The induced map to \(\mathbf A^1_k\) then sends \(x\) to the generic point. Moreover, an irreducible component of \(U\) that meets the generic fibre dominates \(\mathbf A^1_k\), and its generic fibre has dimension one less than the component. Hence the geometric generic fibre has dimension at most \(d-1\), and the induction hypothesis applies to it. Let \(g\) be the projection from the geometric generic fibre of \(U\). By Lemma \hyperref[bp22]{10.7}, locally on \(U\) the complex \(g^*Q\) is isomorphic to the cone of \(\delta_{g^*K}\). This cone vanishes by the induction hypothesis over \(\overline{k(t)}\). As \(x\) lies over the generic point of \(\mathbf A^1_k\), every geometric point above \(x\) factors through the geometric generic fibre. Therefore \(Q\) has zero geometric stalks above every nonclosed point of \(X\).

Only finitely many closed points are left. Indeed, the support of each cohomology sheaf of \(Q\) is constructible, and by the previous paragraph it consists of closed points. A constructible subset of a Noetherian scheme of finite type which contains only closed points is finite, since each positive-dimensional locally closed piece would contain a nonclosed generic point. Thus the finitely many cohomology sheaves of \(Q\) are supported on a single finite set of closed points.

The stalks at these closed points vanish because \(X\) is proper and \(Q\) has no global cohomology. To see that \(R\Gamma(X,Q)=0\), we apply the exchange for proper maps in Lemma \hyperref[bp10]{10.3} to \(a_X\). It identifies \(R\Gamma(X,\delta_K)\) with the double-dual morphism of the perfect complex \(R\Gamma(X,K)\), which is an isomorphism by Lemma \hyperref[bp5]{10.1}. On the other hand, a sheaf \(\mathcal F\) supported on a finite set \(S\) of closed points is \(\bigoplus_{x\in S}x_*\mathcal F_{\bar x}\). Since each residue field is algebraically closed, \(H^0(X,\mathcal F)=\bigoplus_x
\mathcal F_{\bar x}\) and the higher cohomology vanishes. The hypercohomology spectral sequence of the bounded complex \(Q\) therefore degenerates, and it shows that

\[
0=H^mR\Gamma(X,Q)
=\bigoplus_x\mathcal H^m(Q)_{\bar x}.
\]

All these stalks vanish. Since the cohomology sheaves of \(Q\) are supported on this finite set, \(Q=0\).

A finite base field reduces to its algebraic closure. Let \(\bar k\) be an algebraic closure of \(k\) and \(g:X_{\bar k}=X\times_k\bar k\to X\) the projection. The three comparisons of Lemma \hyperref[bp19]{10.6} hold for \(g\), with the same proofs. Indeed, about \(\bar\eta\to B\) these proofs use only that it is a limit of étale \(B\)-schemes, and here, with \(B=\operatorname{Spec}k\) in place of a curve, \(\operatorname{Spec}\bar k\) is the limit of the finite étale \(k\)-schemes \(\operatorname{Spec}k'\), for \(k'\subset\bar k\) finite over \(k\). Unlike in Lemma \hyperref[bp22]{10.7}, the dualizing complex does not change. For an affine open \(V\) of \(X\), choose a closed immersion \(i:V\hookrightarrow\mathbf A^r_k\). By composition of exceptional inverse images and Theorem \bpout{9.3}, \(\omega_V\simeq i^!A(r)[2r]\). The base change \(i_{\bar k}:V_{\bar k}\hookrightarrow\mathbf A^r_{\bar k}\) is a closed immersion into an affine space of the same dimension, so \(\omega_{V_{\bar k}}\simeq i_{\bar k}^!A(r)[2r]\), and the second comparison gives \(g_V^*\omega_V\simeq\omega_{V_{\bar k}}\), with no twist and no shift. The argument of Lemma \hyperref[bp22]{10.7}, with this isomorphism in place of \(g_V^*\omega_V\simeq\omega_{V_{\bar\eta}}(1)[2]\), shows that the isomorphism \(g_V^*\mathbb D_V\mathbb D_V(K|_V)\simeq\mathbb D_{V_{\bar k}}\mathbb D_{V_{\bar k}}g_V^*(K|_V)\) it gives carries \(g_V^*\delta_{K|_V}\) to \(\delta_{g_V^*(K|_V)}\). The latter is an isomorphism by the case already proved, so \(g^*Q\) vanishes on every \(V_{\bar k}\), and \(g^*Q=0\). Finally, every geometric point of \(X\) factors through \(X_{\bar k}\), since its algebraically closed residue field contains an algebraic closure of \(k\). Hence all geometric stalks of \(Q\) vanish, and \(Q=0\) over the finite field as well.

For coefficient rings that are not finite we reduce to finite ones. Indeed, the double-dual map \(\delta:\mathrm{id}\to\mathbb D_X\mathbb D_X\) is a natural transformation of exact functors, built from \(a_X^!\), internal Hom and evaluation. Its extensions of scalars are identified by Proposition \bpout{7.10}. By Proposition \bpout{7.12} the finite case therefore implies the theorem for every admissible \(\Lambda\).
\end{proof}

% text: prose-10-10
\phantomsection\label{bp27}%
Since \(\mathbb D_X\) is now an anti-equivalence, we can dualize the formulas of Lemma \hyperref[bp10]{10.3}. This gives the two exchange formulas with the roles of \(!\) and \(*\) reversed.

% object: corollary-verdier-exchange-formulas (blueprint/chap-Verdier-and-Poincare-duality/sec-Verdier-biduality/corollary-verdier-exchange-formulas/_section.md)
\phantomsection\label{bp28}%
\begin{bpstatementit}{Corollary~10.9}{ (Verdier exchange formulas)}
For a morphism \(f:X\to Y\) there are natural isomorphisms of functors on \(D^b_{\mathrm{ctf}}(X,\Lambda)\), respectively \(D^b_{\mathrm{ctf}}(Y,\Lambda)\),

\[
\begin{aligned}
\mathbb D_YRf_!&\simeq Rf_*\mathbb D_X,&
\mathbb D_YRf_*&\simeq Rf_!\mathbb D_X,\\
\mathbb D_Xf^*&\simeq f^!\mathbb D_Y,&
\mathbb D_Xf^!&\simeq f^*\mathbb D_Y.
\end{aligned}
\]

In particular, for an open \(j:U\hookrightarrow X\) and a closed \(i:Z\hookrightarrow X\),

\[
\mathbb D_Xj_!\simeq Rj_*\mathbb D_U,\qquad
\mathbb D_XRj_*\simeq j_!\mathbb D_U,\qquad
i^*\mathbb D_X\simeq\mathbb D_Zi^!.
\]

The isomorphisms in the left column are (10.3) and (10.4); those in the right column are obtained from them by biduality, as displayed in the proof.
\end{bpstatementit}

% proof: proof/corollary-verdier-exchange-formulas
\begin{proof}[{Proof}]\phantomsection\label{bp29}%
Two of the formulas are Lemma \hyperref[bp10]{10.3}, and we obtain the other two by applying these to the Verdier dual of the argument and using biduality on both schemes. Explicitly, for \(K\in D^b_{\mathrm{ctf}}(X,\Lambda)\) and \(L\in D^b_{\mathrm{ctf}}(Y,\Lambda)\) this reads

\[
\mathbb D_YRf_*K\simeq\mathbb D_YRf_*\mathbb D_X(\mathbb D_XK)
\simeq\mathbb D_Y\mathbb D_YRf_!\mathbb D_XK\simeq Rf_!\mathbb D_XK,
\]

\[
\mathbb D_Xf^!L\simeq\mathbb D_Xf^!\mathbb D_Y(\mathbb D_YL)
\simeq\mathbb D_X\mathbb D_Xf^*\mathbb D_YL\simeq f^*\mathbb D_YL.
\]

The formulas for \(j\) and \(i\) are the special cases \(f=j\) and \(f=i\).
\end{proof}


% object: sec-Poincare-duality (blueprint/chap-Verdier-and-Poincare-duality/sec-Poincare-duality/_section.md), content hidden
\subsection*{\texorpdfstring{10.3\quad Poincaré duality}{10.3 Poincaré duality}}
\phantomsection\label{bp30}%
\addcontentsline{toc}{subsection}{\texorpdfstring{10.3\quad Poincaré duality}{10.3 Poincaré duality}}
% text: prose-10-11
\phantomsection\label{bp31}%
Over an algebraically closed field, \(D^b_{\mathrm{ctf}}(s,\Lambda)\) is the category of perfect \(\Lambda\)-complexes, and on it \(\mathbb D_s\) is \(R\!\operatorname{Hom}_\Lambda(-,\Lambda)\). For \(f=a_X\) the first two Verdier exchange formulas therefore become dualities between compactly supported and ordinary cohomology. This is global Verdier duality (Theorem \hyperref[bp32]{10.10}). It holds for every \(X\) of finite type, smooth or not.

For smooth \(X\) of pure dimension \(d\) we have \(\omega_X\simeq\Lambda_X(d)[2d]\). For locally perfect \(K\) the Verdier dual is therefore the ordinary dual \(K^\vee\) twisted by \((d)[2d]\), and the counit of \(Ra_{X!}\dashv a_X^!\) becomes the trace map of Definition \bpout{9.24}. In this way global duality becomes cup product followed by the trace.

% object: theorem-global-verdier-duality (blueprint/chap-Verdier-and-Poincare-duality/sec-Poincare-duality/theorem-global-verdier-duality/_section.md)
\phantomsection\label{bp32}%
\begin{bpstatementit}{Theorem~10.10}{ (Global Verdier duality)}
Suppose \(k\) is algebraically closed. For every \(K\in D^b_{\mathrm{ctf}}(X,\Lambda)\) there are natural isomorphisms

\begin{align*}
R\Gamma(X,\mathbb D_XK)&\simeq
R\!\operatorname{Hom}_\Lambda(R\Gamma_c(X,K),\Lambda),
\tag{10.5}\\
R\Gamma_c(X,\mathbb D_XK)&\simeq
R\!\operatorname{Hom}_\Lambda(R\Gamma(X,K),\Lambda).
\tag{10.6}\end{align*}

The first is the adjoint of the pairing formed by the cup product with mixed supports, evaluation, and the counit \(\varepsilon_{a_X}\) of \(Ra_{X!}\dashv a_X^!\):

\[
R\Gamma_c(X,K)\otimes^{\mathbf L}R\Gamma(X,\mathbb D_XK)
\longrightarrow R\Gamma_c(X,\omega_X)
\xrightarrow{\varepsilon_{a_X}}\Lambda.
\]
\end{bpstatementit}

% proof: proof/theorem-global-verdier-duality
\begin{proof}[{Proof}]\phantomsection\label{bp33}%
Both isomorphisms are the first two Verdier exchange formulas of Corollary \hyperref[bp28]{10.9} for \(f=a_X\), so we only have to recognize the Verdier dual on the point as the linear dual. On \(s\) the dualizing complex is \(\Lambda_s\), so \(\mathbb D_s=\cHom_s(-,\Lambda_s)\). Since \(k\) is algebraically closed, Lemma \hyperref[bp5]{10.1} identifies \(D^b_{\mathrm{ctf}}(s,\Lambda)\) with \(\operatorname{Perf}(\Lambda)\). Under this identification, then, \(\mathbb D_s\) becomes \(R\!\operatorname{Hom}_\Lambda(-,\Lambda)\), and the two exchange formulas read (10.5) and (10.6). As for the pairing, Lemma \hyperref[bp10]{10.3} constructs the exchange formula (10.3) from evaluation and the counit of \(Rf_!\dashv f^!\). For \(f=a_X\) it is therefore the adjoint of the displayed pairing.
\end{proof}

% text: prose-10-12
\phantomsection\label{bp34}%
Over a finite base field \(k\) the same argument gives us isomorphisms \(Ra_{X*}\mathbb D_XK\simeq\mathbb D_sRa_{X!}K\) and \(Ra_{X!}\mathbb D_XK\simeq\mathbb D_sRa_{X*}K\) in \(D(s_{\mathrm{pro}\acute{e}t},\Lambda)\), with \(\mathbb D_s=\cHom_s(-,\Lambda_s)\). These are the first two formulas of Corollary \hyperref[bp28]{10.9} for \(f=a_X\). After base change to an algebraic closure they become the isomorphisms of the theorem.

The first isomorphism of the theorem is the adjoint of a pairing. We give a name to pairings whose adjoint is an isomorphism.

% object: definition-perfect-pairing (blueprint/chap-Verdier-and-Poincare-duality/sec-Poincare-duality/definition-perfect-pairing/_section.md)
\phantomsection\label{bp35}%
\begin{bpstatement}{Definition~10.11}{ (Perfect pairing)}
A pairing \(P\otimes^{\mathbf L}Q\to\Lambda\) of perfect complexes is \emph{perfect} if its adjoint \(Q\to R\!\operatorname{Hom}_\Lambda(P,\Lambda)\) is an isomorphism.
\end{bpstatement}

% text: prose-10-13
\phantomsection\label{bp36}%
For smooth \(X\), Corollary \bpout{9.28} identifies \(\omega_X\) with \(\Lambda_X(d)[2d]\), and global duality becomes the Poincaré pairing of Weil's axiom (c).

% object: theorem-poincare-duality (blueprint/chap-Verdier-and-Poincare-duality/sec-Poincare-duality/theorem-poincare-duality/_section.md)
\phantomsection\label{bp37}%
\begin{bpstatementit}{Theorem~10.12}{ (Poincaré duality)}
Let \(X/k\) be smooth, separated of pure dimension \(d\), with \(k\) algebraically closed. For every admissible \(\Lambda\) and \(K\in D^b_{\mathrm{ctf}}(X,\Lambda)\), Theorem \bpout{9.3} gives isomorphisms

\[
\omega_X\simeq\Lambda_X(d)[2d],\qquad
\mathbb D_XK\simeq
\cHom_X(K,\Lambda_X)(d)[2d].
\]

Put \(K^\vee=\cHom_X(K,\Lambda_X)\). If \(K\) is locally perfect over \(\Lambda_X\), then so is \(K^\vee\), and it is the ordinary dual. Let \(\operatorname{Tr}_X\) be the map induced on \(R\Gamma_c\) by the trace map of \(a_X\) (Definition \bpout{9.24}). The pairing

\[
\begin{split}
\beta_K:\quad
R\Gamma_c(X,K)\otimes^{\mathbf L}
R\Gamma(X,K^\vee(d))[2d]
&\longrightarrow R\Gamma_c(X,\Lambda_X(d))[2d]\\
&\xrightarrow{\operatorname{Tr}_X}\Lambda
\end{split}
\tag{10.7}
\]

is perfect, and its adjoint is the canonical isomorphism

\[
R\Gamma(X,K^\vee(d))[2d]
\xrightarrow{\sim}
R\!\operatorname{Hom}_\Lambda(R\Gamma_c(X,K),\Lambda).
\tag{10.8}
\]

If \(X\) is proper, \(R\Gamma_c\) may be replaced by \(R\Gamma\). For a map \(\Lambda\to\Lambda'\) of admissible coefficient rings, the change-of-coefficients isomorphisms of Proposition \bpout{7.10} for \(R\Gamma_c\), \(R\Gamma\) and \(\cHom\) carry \(\beta_K\otimes^{\mathbf L}_\Lambda\Lambda'\) to \(\beta_{K_{\Lambda'}}\), where \(K_{\Lambda'}=K\otimes^{\mathbf L}_\Lambda\Lambda'\).
\end{bpstatementit}

% proof: proof/theorem-poincare-duality
\begin{proof}[{Proof}]\phantomsection\label{bp38}%
The theorem is global duality rewritten by means of smooth purity: once \(\omega_X\) and \(\mathbb D_XK\) are identified, the pairing of Theorem \hyperref[bp32]{10.10} becomes \(\beta_K\).

Corollary \bpout{9.28} identifies \(\omega_X\). Since internal Hom commutes with tensoring its second variable by an invertible Tate twist and a shift, we get the displayed description of \(\mathbb D_XK\). If \(K\) is locally perfect, then locally \(K\) is represented by a bounded complex of finite projective modules, whose dual is the termwise dual, with the usual evaluation.

Next, \(\beta_K\) is the pairing of Theorem \hyperref[bp32]{10.10} for \(K\), transported along the isomorphism \(\omega_X\simeq\Lambda_X(d)[2d]\). Indeed, under this isomorphism evaluation \(K\otimes^{\mathbf L}\mathbb D_XK\to\omega_X\) becomes evaluation on coefficients \(K\otimes^{\mathbf L}K^\vee(d)[2d]\to\Lambda_X(d)[2d]\). Moreover, the trace map of Definition \bpout{9.24} is by definition the counit \(\varepsilon_{a_X}\) precomposed with \(\gamma_{a_X}\), so the counit becomes \(\operatorname{Tr}_X\). Hence the adjoint of \(\beta_K\) is the isomorphism (10.5), which is (10.8). By Lemma \hyperref[bp5]{10.1} both factors of \(\beta_K\) are perfect complexes. Since the adjoint of \(\beta_K\) is an isomorphism, \(\beta_K\) is perfect.

For proper \(X\) we have \(Ra_{X!}=Ra_{X*}\), which gives the assertion on proper \(X\). The last statement holds because the pairing is built from tensor products, the two direct images, evaluation, the isomorphism \(\gamma_{a_X}\) and the trace map, all of which Proposition \bpout{7.10} and Theorem \bpout{9.3} show to be compatible with every coefficient map.
\end{proof}

% text: prose-10-14
\phantomsection\label{bp39}%
For \(X=\mathbf A^1\) and \(K=\Lambda\) the theorem pairs \(R\Gamma_c(\mathbf A^1,\Lambda)=\Lambda(-1)[-2]\) with \(R\Gamma(\mathbf A^1,\Lambda(1))[2]=\Lambda(1)[2]\), and the pairing is the trace \(H^2_c(\mathbf A^1,\Lambda(1))\simeq\Lambda\). Over a field we can pass from this statement about complexes to statements about the individual cohomology groups.

% object: corollary-degreewise-duality-over-fields (blueprint/chap-Verdier-and-Poincare-duality/sec-Poincare-duality/corollary-degreewise-duality-over-fields/_section.md)
\phantomsection\label{bp40}%
\begin{bpstatementit}{Corollary~10.13}{ (Degreewise duality over fields)}
Let \(X/k\) be smooth and separated of pure dimension \(d\), with \(k\) algebraically closed, let \(\Lambda\) be a field (a finite coefficient field, \(\mathbf Q_\ell\), \(E\) or \(\overline{\mathbf Q}_\ell\)), let \(K\in D^b_{\mathrm{ctf}}(X,\Lambda)\), and put \(K^\vee=\cHom_X(K,\Lambda_X)\). For every \(i\) the pairings induced by \(\beta_K\),

\[
H_c^i(X,K)\times H^{2d-i}(X,K^\vee(d))\longrightarrow\Lambda,
\qquad
H_c^i(X,K)\times H^{2d-i}(X,K^\vee)\longrightarrow\Lambda(-d),
\]

are perfect pairings of finite-dimensional vector spaces. If \(X\) is proper, the same holds with \(H_c^i\) replaced by \(H^i\).
\end{bpstatementit}

% proof: proof/corollary-degreewise-duality-over-fields
\begin{proof}[{Proof}]\phantomsection\label{bp41}%
We take cohomology of the isomorphism (10.8); over a field this loses nothing, because linear duality is exact. More precisely, a bounded complex with finite-dimensional cohomology is isomorphic in the derived category to the sum of its shifted cohomology spaces, so that \(H^{-i}(R\!\operatorname{Hom}_\Lambda(P,\Lambda))=\Hom_\Lambda(H^i(P),
\Lambda)\). Because of the shift \([2d]\), the source of (10.8) has \(H^{2d-i}(X,K^\vee(d))\) in degree \(-i\), and taking \(H^{-i}\) we get \(H^{2d-i}(X,K^\vee(d))\simeq\Hom_\Lambda(H^i_c(X,K),\Lambda)\). Since this isomorphism is induced by the pairing in degrees \((i,2d-i)\), it says exactly that the first pairing is perfect. Here both spaces are finite-dimensional by Lemma \hyperref[bp5]{10.1}. Finally, the second pairing is the first one with the second factor and the value line twisted by \((-d)\), so it is perfect as well.
\end{proof}

% text: prose-10-15
\phantomsection\label{bp42}%
The only property of a field used in this proof is that \(R\!\operatorname{Hom}_\Lambda(-,\Lambda)\) commutes with taking cohomology.

% object: remark-duality-cohomology-groups-over-oe (blueprint/chap-Verdier-and-Poincare-duality/sec-Poincare-duality/remark-duality-cohomology-groups-over-oe/_section.md)
\phantomsection\label{bp43}%
\begin{bpstatement}{Remark~10.14}{ (Duality on cohomology groups over $\mathcal O_E$ and finite rings)}
Over \(\mathcal O_E\), derived duality gives only the natural universal coefficient exact sequence

\[
0\longrightarrow
\Ext^1_{\mathcal O_E}(H_c^{i+1}(X,K),\mathcal O_E)
\longrightarrow H^{2d-i}(X,K^\vee(d))
\longrightarrow
\Hom_{\mathcal O_E}(H_c^i(X,K),\mathcal O_E)
\longrightarrow0.
\]

It comes from applying derived Hom to a bounded finite free model of \(R\Gamma_c(X,K)\), since \(\mathcal O_E\) has global dimension one. The sequence need not split canonically, and because of its \(\Ext^1\) term the pairings on individual cohomology modules need not be perfect. For example, on a point the perfect complex \([\mathcal O_E\xrightarrow{\pi_E}\mathcal O_E]\) in degrees \(-1,0\) has \(H^0=\mathcal O_E/\pi_E\), whereas its dual has \(H^0=0\) and \(H^1=\mathcal O_E/\pi_E\). Derived evaluation is an isomorphism, but the degree-zero pairing is not perfect.

Over a finite ring \(A\) that is not injective as a module over itself, derived Hom out of a bounded complex of finite projective modules records all the groups \(\Ext^i_A(-,A)\) of the cohomology modules, and this information may be lost in their ordinary duals. Moreover, the cohomology modules of a complex in \(D^b_{\mathrm{ctf}}(X,A)\) need not have finite Tor dimension, so the derived pairing cannot be reduced to pairings of cohomology sheaves.
\end{bpstatement}


% object: sec-Cohomology-in-top-degree (blueprint/chap-Verdier-and-Poincare-duality/sec-Cohomology-in-top-degree/_section.md), content hidden
\subsection*{\texorpdfstring{10.4\quad Cohomology in top degree}{10.4 Cohomology in top degree}}
\phantomsection\label{bp44}%
\addcontentsline{toc}{subsection}{\texorpdfstring{10.4\quad Cohomology in top degree}{10.4 Cohomology in top degree}}
% text: prose-10-16
\phantomsection\label{bp45}%
For constant coefficients on a smooth, proper and connected \(X\), we expect \(H^{2d}(X,\Lambda_X(d))\) to be free of rank one, generated by the class of a point, as for a manifold. We use this normalization in the trace formula and in point counting.

For a compact connected oriented manifold this is the statement that integration \(H^n(M)\to\mathbf R\) is an isomorphism sending a point to \(1\). Here the statement also has an arithmetic meaning, because geometric Frobenius acts on \(\Lambda(-d)\) by \(q^d\). Once we know that the trace map commutes with Frobenius (Theorem \hyperref[bp65]{10.21}), the theorem below says that Frobenius acts on \(H^{2d}\) by \(q^d\) and on \(H^0\) trivially. These give the factors \(1-T\) and \(1-q^dT\) of the zeta function.

Combined with Corollary \hyperref[bp40]{10.13}, the theorem gives the cup product pairing (Corollary \hyperref[bp52]{10.17}).

% object: theorem-trace-isomorphism-top-degree (blueprint/chap-Verdier-and-Poincare-duality/sec-Cohomology-in-top-degree/theorem-trace-isomorphism-top-degree/_section.md)
\phantomsection\label{bp46}%
\begin{bpstatementit}{Theorem~10.15}{ (The trace isomorphism in top degree)}
Let \(X/k\) be smooth, proper and connected of pure dimension \(d\), with \(k\) algebraically closed. For every admissible coefficient ring,

\[
H^i(X,\Lambda_X)=0\quad(i\notin[0,2d]),\qquad
H^0(X,\Lambda_X)=\Lambda,
\]

and the map

\[
\operatorname{tr}_X:H^{2d}(X,\Lambda_X)
\xrightarrow{\sim}\Lambda(-d),
\]

the twist by \((-d)\) of the map \(\operatorname{tr}_X(d):H^{2d}(X,\Lambda_X(d))\to\Lambda\) induced on \(H^{2d}\) by \(\operatorname{Tr}_X\) of Definition \bpout{9.24}, is an isomorphism.
\end{bpstatementit}

% proof: proof/theorem-trace-isomorphism-top-degree
\begin{proof}[{Proof}]\phantomsection\label{bp47}%
We read off the top degree from Poincaré duality in degree zero over \(\mathbf Z/\ell^n\), and then change coefficients: to a finite \(A\) by extension of scalars, to \(\mathcal O_E\) by a limit, and to the fields by inverting \(\ell\).

With finite coefficients, étale cohomology of a \(d\)-dimensional scheme of finite type over an algebraically closed field vanishes above degree \(2d\), that is, \(\operatorname{cd}_\ell(X)\le2d\) in the notation of Definition \bpout{2.45} [\hyperref[bp70]{SGA4}, X, Cor. 4.3]. This gives the vanishing outside \([0,2d]\), and \(H^0=\Lambda\) because \(X\) is connected.

Now let \(\Lambda=\mathbf Z/\ell^n\). Since this ring is injective as a module over itself, \(R\!\operatorname{Hom}_{\mathbf Z/\ell^n}(-,\mathbf Z/\ell^n)\) has cohomology equal to the ordinary dual of the opposite cohomology. Hence Poincaré duality (Theorem \hyperref[bp37]{10.12}) in degree zero identifies \(H^{2d}(X,\mathbf Z/\ell^n(d))\) with \(\Hom(H^0(X,\mathbf Z/\ell^n),\mathbf Z/\ell^n)=\mathbf Z/\ell^n\). By Theorem \hyperref[bp37]{10.12} this identification is given by pairing with \(H^0\) and applying the trace map, so it is \(\operatorname{tr}_X(d)\).

A general finite ring \(A\) need not be injective over itself, so we reach it by extension of scalars instead. Every finite \(A\) is an algebra over some \(\mathbf Z/\ell^n\), and by Proposition \bpout{7.10} the complexes with constant coefficients and their traces over \(A\) are obtained by derived extension of scalars from \(\mathbf Z/\ell^n\). For a complex \(P\) with cohomology in degrees \(\le2d\), right exactness of the tensor product gives \(H^{2d}(P\otimes^{\mathbf L}A)=H^{2d}(P)\otimes A\). The top cohomology is therefore \(\mathbf Z/\ell^n\otimes A=A\), and the trace is an isomorphism. Since derived extension of scalars does not raise the top degree, the same upper bound holds.

For \(\mathcal O_E\), we write \(R\Gamma(X,\mathcal O_{E,X}(d))\) as the derived inverse limit of the complexes \(R\Gamma(X,(\mathcal O_E/\pi_E^n)_X(d))\). Its cohomology fits into the Milnor sequence \(0\to\lim^1_nH^{i-1}\to H^i\to\lim_nH^i\to0\). Since the cohomology groups at each level are finite, they satisfy the Mittag–Leffler condition and there is no \(\lim^1\) contribution to the top group. The traces in top degree are compatible, so we get \(H^{2d}(X,\mathcal O_{E,X}(d))=\mathcal O_E\), with the same vanishing bounds. The lower bound and \(H^0\) also follow directly from ordinary sheaf cohomology and connectedness. Finally, inverting \(\ell\) and then extending scalars along finite extensions of \(E\), we obtain the cases \(\mathbf Q_\ell\), \(E\) and \(\overline{\mathbf Q}_\ell\).
\end{proof}

% text: prose-10-17
\phantomsection\label{bp48}%
By the theorem, \(H^{2d}(X,\Lambda_X(d))\) is free of rank one. It is generated by the class of a point, which we define by the counit of the exceptional adjunction for the inclusion of the point. Let \(x:\operatorname{Spec}k\to X\) be a closed point. Using \(x^!\omega_X=(a_Xx)^!\Lambda_s=\Lambda\) and \(\omega_X\simeq\Lambda_X(d)[2d]\), the counit \(x_*x^!\omega_X\to\omega_X\) induces a map \(\Lambda\to R\Gamma(X,\omega_X)=R\Gamma(X,\Lambda_X(d))[2d]\), and we let \([x]\in H^{2d}(X,\Lambda_X(d))\) be the image of \(1\).

% object: lemma-class-of-point (blueprint/chap-Verdier-and-Poincare-duality/sec-Cohomology-in-top-degree/lemma-class-of-point/_section.md)
\phantomsection\label{bp49}%
\begin{bpstatementit}{Lemma~10.16}{ (The class of a point)}
In the setting of Theorem \hyperref[bp46]{10.15}, we have \(\operatorname{tr}_X(d)([x])=1\) for every closed point \(x\) of \(X\). In particular \([x]\) generates \(H^{2d}(X,\Lambda_X(d))\).
\end{bpstatementit}

% proof: proof/lemma-class-of-point
\begin{proof}[{Proof}]\phantomsection\label{bp50}%
By definition, \([x]\) is the image of \(1\) under the map induced by the counit \(x_*x^!\omega_X\to\omega_X\), and \(\operatorname{tr}_X(d)\) is induced by the counit \(Ra_{X*}\omega_X\to\Lambda\). Because exceptional adjunctions compose, the composite of these two counits is the counit for \(a_Xx=\mathrm{id}_s\), which is the identity of \(\Lambda\). Thus \(\operatorname{tr}_X(d)([x])=1\). Since \(\operatorname{tr}_X(d)\) is an isomorphism onto \(\Lambda\) by Theorem \hyperref[bp46]{10.15}, \([x]\) is a generator.
\end{proof}

% text: prose-10-18
\phantomsection\label{bp51}%
For a field \(\Lambda\) and the constant sheaf, the trace isomorphism in top degree turns Poincaré duality into the classical cup product pairing. On \(\mathbf P^1\), for instance, the only nontrivial instance is \(H^0\times H^2\to\Lambda(-1)\). After twisting by \((1)\) it pairs the generators \(1\) and \([x]\) to \(1\) by Lemma \hyperref[bp49]{10.16}, so it is perfect.

% object: corollary-cup-product-pairing (blueprint/chap-Verdier-and-Poincare-duality/sec-Cohomology-in-top-degree/corollary-cup-product-pairing/_section.md)
\phantomsection\label{bp52}%
\begin{bpstatementit}{Corollary~10.17}{ (The cup product pairing)}
Let \(X/k\) be smooth, proper and connected of pure dimension \(d\), with \(k\) algebraically closed, and let \(\Lambda\) be a field. For every \(i\) the pairing

\[
\langle a,b\rangle=\operatorname{tr}_X(a\cup b),\qquad
H^i(X,\Lambda)\times H^{2d-i}(X,\Lambda)\longrightarrow\Lambda(-d),
\]

is perfect, and \(\langle b,a\rangle=(-1)^i\langle a,b\rangle\). In particular, if \(d\) is odd and \(\Lambda\) has characteristic zero, the pairing on \(H^d(X,\Lambda)\) is alternating, as for a compact complex manifold of odd dimension.
\end{bpstatementit}

% proof: proof/corollary-cup-product-pairing
\begin{proof}[{Proof}]\phantomsection\label{bp53}%
For \(K=\Lambda_X\) we have \(K^\vee=\Lambda_X\), and since \(X\) is proper, Corollary \hyperref[bp40]{10.13} says that the pairing \(H^i\times H^{2d-i}\to\Lambda(-d)\) induced by \(\beta_{\Lambda_X}\) is perfect. This pairing is \(\langle-,-\rangle\), because it is cup product followed by the trace map in top degree, that is, by the isomorphism \(\operatorname{tr}_X\) of Theorem \hyperref[bp46]{10.15}. By graded commutativity of the cup product we have \(\langle b,a\rangle=(-1)^{i(2d-i)}\langle a,b\rangle
=(-1)^i\langle a,b\rangle\).
\end{proof}


% object: sec-Weil-complexes (blueprint/chap-Verdier-and-Poincare-duality/sec-Weil-complexes/_section.md), content hidden
\subsection*{\texorpdfstring{10.5\quad Weil complexes}{10.5 Weil complexes}}
\phantomsection\label{bp54}%
\addcontentsline{toc}{subsection}{\texorpdfstring{10.5\quad Weil complexes}{10.5 Weil complexes}}
% text: prose-10-19
\phantomsection\label{bp55}%
The Weil conjectures concern varieties over \(\mathbf F_q\), while the six operations of Part I live on schemes over the algebraic closure. On \(X=X_0\times_{\mathbf F_q}k\), what remains of a variety \(X_0\) over \(\mathbf F_q\) is an action of Frobenius on \(X\) and on every complex pulled back from \(X_0\).

Here \(k_0=\mathbf F_q\) and \(k\) is an algebraic closure of \(k_0\). In this setting the geometric Frobenius \(\sigma\in\operatorname{Gal}(k/k_0)\) is the inverse of the arithmetic Frobenius \(a\mapsto a^q\). For a scheme \(T\) of the form \(T_0\times_{k_0}k\), with \(T_0\) over \(k_0\), let

\[
\operatorname{Fr}_T:T\longrightarrow T
\]

be the semilinear automorphism induced by \(\sigma\). It is a morphism of schemes covering the automorphism \(\operatorname{Spec}(\sigma)\) of \(\operatorname{Spec}k\).

We use \(\operatorname{Fr}_T\) to define Weil structures, but we prove the trace formula for the \(k\)-linear Frobenius \(\Phi_T=F_{T_0/k_0}\times_{k_0}k\), obtained by base change from the relative \(q\)-power Frobenius of \(T_0\), since its graph is a correspondence over \(k\). The two are compared in Section \bpout{12.1} (Lemma \bpout{12.3}).

Since \(\operatorname{Fr}_T\) is an automorphism of schemes, pullback of sheaves of modules along it is exact and preserves derived completeness. Hence for every admissible coefficient ring \(\Lambda\) it induces an autoequivalence

\[
\operatorname{Fr}_T^*:
D^b_{\mathrm{ctf}}(T,\Lambda)\longrightarrow
D^b_{\mathrm{ctf}}(T,\Lambda),
\]

which agrees with the derived inverse image along \(\operatorname{Fr}_T\).

Instead of complexes on \(T_0\) pulled back to \(T\), we keep only the part of that datum that the trace formula sees: the complex on \(T\), where all cohomology is computed, together with an identification of it with its Frobenius pullback. Every construction on \(T\) (direct images, \(f^!\), truncations, duals) is compatible with such identifications by base change, as we check below, so we never have to descend anything to \(T_0\).

There are also more of these pairs than complexes on \(T_0\), and we use the extra ones in Part III. Indeed, for a \(\mathbf Q_\ell\)-sheaf on \(T_0\), the stalk at a geometric point over a closed point is a continuous representation of a Galois group. It follows that Frobenius preserves a lattice in it, and its eigenvalues are \(\ell\)-adic units. By contrast, on \(\operatorname{Spec}k\) the sheaf \(\mathbf Q_\ell\) with Frobenius acting by \(\ell\) is a pair of this kind that does not come from \(\operatorname{Spec}\mathbf F_q\). Pulled back to a scheme, such pairs are the sheaves \(\mathcal B_b\) that normalize weights in Proposition \bpout{13.7}.

We require an isomorphism so that the Frobenius operator on cohomology is invertible over a field (Theorem \hyperref[bp65]{10.21}). The trace formula itself only needs a morphism (Definition \bpout{12.2}), and it is stated for these pairs (Theorem \bpout{12.24}).

% object: definition-weil-structure (blueprint/chap-Verdier-and-Poincare-duality/sec-Weil-complexes/definition-weil-structure/_section.md)
\phantomsection\label{bp56}%
\begin{bpstatement}{Definition~10.18}{ (Weil structure)}
For \(K\in D^b_{\mathrm{ctf}}(T,\Lambda)\), a \emph{Weil structure} is an isomorphism

\[
\phi_K:\operatorname{Fr}_T^*K\xrightarrow{\sim}K.
\]

A \emph{Weil complex} is a pair \((K,\phi_K)\) with this structure. A morphism \((K,\phi_K)\to(L,\phi_L)\) is a morphism \(h:K\to L\) satisfying

\[
h\circ\phi_K=\phi_L\circ\operatorname{Fr}_T^*(h).
\]
\end{bpstatement}

% text: prose-10-20
\phantomsection\label{bp57}%
A Weil structure is the datum needed to define a Frobenius operator on cohomology. The basic example is a complex \(K\) pulled back from a complex \(K_0\) on \(T_0\). Since the projection \(T\to T_0\) composed with \(\operatorname{Fr}_T\) is again the projection, the pullback of \(K_0\) gives a canonical isomorphism \(\operatorname{Fr}_T^*K\simeq K\).

The triangles, truncations and spectral sequences we form in \(D^b_{\mathrm{ctf}}(T,\Lambda)\) are compatible with Weil structures. Since \(\operatorname{Fr}_T^*\) is exact on sheaves, it commutes with standard truncations and with the triangle of a short exact sequence of sheaves. It also commutes with the localization triangle \(j_!j^*K\to K\to i_*i^*K\) for an open immersion \(j:U\to T\) obtained by base change from an open \(U_0\subset T_0\), with closed complement \(i\). Indeed, for such \(U\) we have \(\operatorname{Fr}_T^{-1}(U)=U\), and \(\operatorname{Fr}_T\) restricts to \(\operatorname{Fr}_U\) and to the Frobenius of the complement. These constructions therefore carry Weil structures to Weil structures, and the induced long exact cohomology sequences are Frobenius-equivariant.

Let \((M,\phi_M)\) be a Weil complex on \(X\) and let \(f:X\to Y\) come from \(f_0\) over \(k_0\). The base change map for the cartesian square with horizontal arrows \(\operatorname{Fr}_X\), \(\operatorname{Fr}_Y\) identifies \(\operatorname{Fr}_Y^*Rf_{!}M\) with \(Rf_{!}\operatorname{Fr}_X^*M\), and the composite

\[
\operatorname{Fr}_Y^*Rf_{!}M
\simeq Rf_{!}\operatorname{Fr}_X^*M
\xrightarrow{\ Rf_{!}(\phi_M)\ }Rf_{!}M
\]

is a Weil structure on \(Rf_!M\). Likewise \(f^*\), \(Rf_*\), \(f^!\), \(\otimes^{\mathbf L}\) and \(\cHom\) carry Weil structures, via the corresponding base change maps. These maps are isomorphisms for every \(f\), because \(\operatorname{Fr}_X\) and \(\operatorname{Fr}_Y\) are isomorphisms of schemes and base change along an isomorphism is transport of structure.

For \(Rf_!\) this needs a word, since \(\operatorname{Fr}_X\) and \(\operatorname{Fr}_Y\) are not \(k\)-morphisms, while \(Rf_!\) is built from compactifications over \(k\). We compute \(Rf_!=R\bar f_*j_!\) with a compactification \(X\overset{j}{\hookrightarrow}\bar X\xrightarrow{\bar f}Y\) obtained by base change from a compactification of \(f_0\) over \(k_0\) (Theorem \bpout{A.2}). Then \(\operatorname{Fr}_{\bar X}\) is an automorphism of the scheme \(\bar X\) with \(\operatorname{Fr}_{\bar X}^{-1}(X)=X\) and \(\bar f\operatorname{Fr}_{\bar X}=\operatorname{Fr}_Y\bar f\). For an automorphism \(h\) of a scheme, \(h^*\simeq(h^{-1})_*\), so transport of structure gives \(\operatorname{Fr}_Y^*R\bar f_*\simeq R\bar f_*\operatorname{Fr}_{\bar X}^*\) and \(\operatorname{Fr}_{\bar X}^*j_!\simeq j_!\operatorname{Fr}_X^*\), and their composite is the base change map for \(Rf_!\). It does not depend on the compactification. Indeed, two compactifications of \(f_0\) are dominated by a third, the closure of \(X_0\) in their fibre product over \(Y_0\), and the refinement isomorphisms of \(Rf_!\) come from proper maps over \(k_0\), which commute with the Frobenii. The map for \(f^!\) is the mate of the map for \(Rf_!\), and it is an isomorphism because both functors \(\operatorname{Fr}^*\) are equivalences.

Pullback along \(\operatorname{Fr}_T\) on compactly supported cohomology is constructed in the same way. For a morphism of schemes \(h:S'\to S\) and \(L\in D(S_{\mathrm{pro}\acute{e}t},\Lambda)\), the \emph{pullback map}

\[
\operatorname{pb}_{h,L}:R\Gamma(S,L)\longrightarrow R\Gamma(S',h^*L)
\]

is \(R\Gamma\) of the unit \(L\to Rh_*h^*L\), followed by the identification \(R\Gamma(S,Rh_*-)=R\Gamma(S',-)\). We write \(\operatorname{pb}_h\) when \(L\) is clear. It is natural in \(L\), and it is transitive: for \(g:S'\to S\) and \(h:S''\to S'\) we have \(\operatorname{pb}_{gh,L}=\operatorname{pb}_{h,g^*L}\circ\operatorname{pb}_{g,L}\) under \((gh)^*\simeq h^*g^*\). Let \(j:T\hookrightarrow\bar T\) be the base change of a compactification of \(T_0\) over \(k_0\). Then \(\operatorname{Fr}_{\bar T}\) is an automorphism of \(\bar T\) with \(\operatorname{Fr}_{\bar T}^{-1}(T)=T\), and \emph{pullback along \(\operatorname{Fr}_T\)} is the composite

\[
\operatorname{pb}_{\operatorname{Fr}_T}:
R\Gamma_c(T,K)=R\Gamma(\bar T,j_!K)
\xrightarrow{\ \operatorname{pb}_{\operatorname{Fr}_{\bar T}}\ }
R\Gamma(\bar T,\operatorname{Fr}_{\bar T}^*j_!K)
\simeq R\Gamma(\bar T,j_!\operatorname{Fr}_T^*K)
=R\Gamma_c(T,\operatorname{Fr}_T^*K).
\]

It does not depend on \(\bar T\). Indeed, a refinement \(\pi:\bar T'\to\bar T\) between two such compactifications comes from \(k_0\), so \(\pi\operatorname{Fr}_{\bar T'}=\operatorname{Fr}_{\bar T}\pi\). Transitivity of pullback maps for this composite, together with the base change \(\operatorname{Fr}_{\bar T}^*R\pi_*\simeq R\pi_*\operatorname{Fr}_{\bar T'}^*\) along the automorphisms \(\operatorname{Fr}\) and with \(R\pi_*j'_!\simeq j_!\), shows that the refinement isomorphism \(R\Gamma(\bar T',j'_!K)\simeq R\Gamma(\bar T,R\pi_*j'_!K)
\simeq R\Gamma(\bar T,j_!K)\) carries one pullback map to the other. In the same way, the pullback map along \(\operatorname{Fr}_T^m\) is defined, and by transitivity it is the \(m\)-fold composite \(\operatorname{pb}_{\operatorname{Fr}_T^m,K}=\operatorname{pb}_{\operatorname{Fr}_T,(\operatorname{Fr}_T^{m-1})^*K}\circ\cdots\circ\operatorname{pb}_{\operatorname{Fr}_T,K}\).

% object: definition-frobenius-operator (blueprint/chap-Verdier-and-Poincare-duality/sec-Weil-complexes/definition-frobenius-operator/_section.md)
\phantomsection\label{bp58}%
\begin{bpstatement}{Definition~10.19}{ (Frobenius operator)}
For a Weil complex \((K,\phi_K)\) on \(T\), the \emph{geometric Frobenius operator} on compactly supported cohomology is the endomorphism

\[
\begin{gathered}
\operatorname{Frob}_{T,K}:R\Gamma_c(T_{\mathrm{pro}\acute{e}t},K)
\xrightarrow{\ \operatorname{pb}_{\operatorname{Fr}_T}\ }
R\Gamma_c(T_{\mathrm{pro}\acute{e}t},\operatorname{Fr}_T^*K)\\
\xrightarrow{\ R\Gamma_c(T_{\mathrm{pro}\acute{e}t},\phi_K)\ }
R\Gamma_c(T_{\mathrm{pro}\acute{e}t},K).
\end{gathered}
\]

Here \(\operatorname{pb}_{\operatorname{Fr}_T}\) is the pullback along \(\operatorname{Fr}_T\) constructed above.
\end{bpstatement}

% text: prose-10-21
\phantomsection\label{bp59}%
When \(T\) and \(K\) are clear, we write \(F\) for \(\operatorname{Frob}_{T,K}\) and for the endomorphisms it induces on the groups \(H^i_c(T,K)\).

For \(r\in\mathbf Z\), write \(\Lambda(r)=\Gamma((\operatorname{Spec}k)_{\mathrm{pro}\acute{e}t},
\Lambda_{\operatorname{Spec}k}(r))\) for the Tate coefficient module with its induced Frobenius operator. Arithmetic Frobenius acts on \(\mu_{\ell^n}(k)\) by \(\zeta\mapsto\zeta^q\). Hence geometric Frobenius acts on \(\Lambda(1)\) by \(q^{-1}\), and on \(\Lambda(-d)\) by \(q^d\). Let \(X_0\) be smooth of pure dimension \(d\). After removing the twist, the trace map of Definition \bpout{9.24} is a Frobenius-equivariant map, by Proposition \hyperref[bp62]{10.20} below:

\[
H_c^{2d}(X_{\mathrm{pro}\acute{e}t},\Lambda_X)\longrightarrow\Lambda(-d).
\]

For \(X=\mathbf P^1_k\) this says that geometric Frobenius acts on \(H^2(\mathbf P^1_k,\Lambda)=\Lambda(-1)\) by \(q\), as checked after Proposition \bpout{6.7}.


% object: sec-Frobenius-and-Poincare-duality (blueprint/chap-Verdier-and-Poincare-duality/sec-Frobenius-and-Poincare-duality/_section.md), content hidden
\subsection*{\texorpdfstring{10.6\quad Frobenius and Poincaré duality}{10.6 Frobenius and Poincaré duality}}
\phantomsection\label{bp60}%
\addcontentsline{toc}{subsection}{\texorpdfstring{10.6\quad Frobenius and Poincaré duality}{10.6 Frobenius and Poincaré duality}}
% text: prose-10-22
\phantomsection\label{bp61}%
Let \(X_0\) be smooth of pure dimension \(d\) and let \(K\) be locally perfect, so that \(\mathbb D_XK\simeq K^\vee(d)[2d]\). In this case cup product pairs \(H^i_c(X,K)\) with \(H^{2d-i}(X,K^\vee)\) and takes values in \(H^{2d}_c(X,\Lambda)\). To use this pairing for the Weil conjectures we need to know how \(F\) acts on both sides. The target is not the trivial Weil complex but \(\Lambda(-d)\), as for \(\mathbf P^1\) above. If we forgot the twist, we would predict the relation \(\alpha\leftrightarrow1/\alpha\) between eigenvalues in complementary degrees instead of \(\alpha\leftrightarrow q^d/\alpha\).

We built Poincaré duality from the purity isomorphism \(\gamma_f\) and the trace map. It therefore suffices to check that these two maps commute with pullback along Frobenius (Proposition \hyperref[bp62]{10.20}). This statement about natural transformations is transport of structure, except for the identification of Tate twists. There we use that pullback along \(\operatorname{Fr}_T\) respects Kummer classes. The compatibility of the duality pairing with Weil structures then follows formally (Theorem \hyperref[bp65]{10.21}).

Let \(f\) be obtained from a smooth morphism \(f_0:X_0\to Y_0\) over \(k_0=\mathbf F_q\) by extension to \(k\), let \(\operatorname{Fr}_X\), \(\operatorname{Fr}_Y\) be the geometric Frobenius maps above, and let \(\gamma_f\) be the isomorphism of (9.2).

% object: proposition-purity-trace-commute-frobenius (blueprint/chap-Verdier-and-Poincare-duality/sec-Frobenius-and-Poincare-duality/proposition-purity-trace-commute-frobenius/_section.md)
\phantomsection\label{bp62}%
\begin{bpstatementit}{Proposition~10.20}{ (Purity and trace commute with Frobenius)}
Fix an admissible \(\Lambda\), and let \(K\in D^b_{\mathrm{ctf}}(Y,\Lambda)\) and \(M\in D^b_{\mathrm{ctf}}(X,\Lambda)\). The square with horizontal arrows \(\operatorname{Fr}_X,\operatorname{Fr}_Y\) and vertical arrows \(f,f\) is cartesian, and its base change comparisons identify

\[
\operatorname{Fr}_X^*f^!K
\simeq f^!\operatorname{Fr}_Y^*K,
\qquad
\operatorname{Fr}_Y^*Rf_!M
\simeq Rf_!\operatorname{Fr}_X^*M.
\tag{10.9}
\]

With these identifications, and the Kummer identification \(\operatorname{Fr}_T^*\Lambda_T(1)\simeq\Lambda_T(1)\) for the twists, we have

\[
\operatorname{Fr}_X^*\gamma_f(K)
=\gamma_f(\operatorname{Fr}_Y^*K),
\qquad
\operatorname{Fr}_Y^*\operatorname{Tr}_f(K)
=\operatorname{Tr}_f(\operatorname{Fr}_Y^*K).
\]
\end{bpstatementit}

% proof: proof/proposition-purity-trace-commute-frobenius
\begin{proof}[{Proof}]\phantomsection\label{bp63}%
Both identities are instances of base change for smooth purity, applied to a square whose horizontal maps are not \(k\)-linear. Since \(\gamma_f\) and \(\operatorname{Tr}_f\) are built from base change maps and from the Kummer class of the twist, we only have to know how \(\operatorname{Fr}_T\) acts on Kummer classes.

The first identity is Proposition \bpout{9.22} for this square. Since \(\operatorname{Fr}_X\) is not a morphism over \(k\), the only point to check is that the identification of twists used there is the Kummer identification, that is, that pullback along \(\operatorname{Fr}_T\) sends the Kummer class of a unit to the Kummer class of its pullback. This holds because the Kummer sequence \(0\to\mu_{\ell^n}\to\mathbf G_m\xrightarrow{\ell^n}\mathbf G_m\to0\) and its boundary map are natural for every morphism of schemes, \(k\)-linear or not.

For the second identity we combine the first with the counit identity defining exceptional base change, which says that after base change for compact support, the pulled-back counit is the counit preceded by the exceptional comparison. As the trace map is the counit composed with \(\gamma_f\), the two identities together give the statement for \(\operatorname{Tr}_f\).
\end{proof}

% text: prose-10-23
\phantomsection\label{bp64}%
Weil structures are transported along \(f^!\), \(f^*\) and \(Rf_!\). On \(f^*K(d)[2d]\) the Weil structure is the pullback of \(\phi_K\) combined with the Kummer identification. By the proposition above and naturality of \(\gamma_f\) and \(\operatorname{Tr}_f\) applied to \(\phi_K\), the maps \(\gamma_f(K)\) and \(\operatorname{Tr}_f(K)\) are morphisms of Weil complexes.

The constant sheaf \(\Lambda_Y\) has the Weil structure given by the identification \(\operatorname{Fr}_Y^*\Lambda_Y\simeq\Lambda_Y\) of pullbacks of constant sheaves. For \(K=\Lambda_Y\) with this structure, the previous paragraph says that the isomorphism \(\omega_f\simeq\Lambda_X(d)[2d]\) and the trace \(\operatorname{Tr}_f(\Lambda_Y)\) are morphisms of Weil complexes.

% object: theorem-poincare-duality-frobenius (blueprint/chap-Verdier-and-Poincare-duality/sec-Frobenius-and-Poincare-duality/theorem-poincare-duality-frobenius/_section.md)
\phantomsection\label{bp65}%
\begin{bpstatementit}{Theorem~10.21}{ (Poincaré duality and Frobenius)}
Let \(X_0/\mathbf F_q\) be smooth and separated of pure dimension \(d\), let \(X=X_0\times\overline{\mathbf F}_q\), and let \((K,\phi_K)\) be a Weil complex on \(X\) (Definition \hyperref[bp56]{10.18}) with \(K\) locally perfect. Equip the dual with the contragredient Weil structure

\[
\operatorname{Fr}_X^*K^\vee
\simeq(\operatorname{Fr}_X^*K)^\vee
\xrightarrow{\ (\phi_K^{-1})^\vee\ }K^\vee.
\]

For every admissible \(\Lambda\) the Poincaré pairing \(\beta_K\) of (10.7) is a morphism of Weil complexes. Thus, with \(F\) the Frobenius operators defined above, the pairing \(H^i_c(X,K)\times H^{2d-i}(X,K^\vee(d))\to\Lambda\) satisfies \(\langle Fx,Fy\rangle=\langle x,y\rangle\), and the pairing \(H^i_c(X,K)\times H^{2d-i}(X,K^\vee)\to\Lambda(-d)\) satisfies

\[
\langle Fx,Fy\rangle=q^d\langle x,y\rangle.
\]

If \(\Lambda\) is a field, then \(F\) is invertible on \(H^i_c(X,K)\), and with respect to the second pairing \(F\) on \(H^{2d-i}(X,K^\vee)\) is \(q^d\) times the inverse transpose of \(F\) on \(H^i_c(X,K)\). Hence the two spaces have the same dimension \(b\), the eigenvalues of \(F\) on \(H^{2d-i}(X,K^\vee)\) are the \(q^d/\alpha\), with multiplicities, where \(\alpha\) ranges over the eigenvalues on \(H^i_c(X,K)\), and \(\det(F\mid H^i_c(X,K))\det(F\mid H^{2d-i}(X,K^\vee))=q^{db}\).
\end{bpstatementit}

% proof: proof/theorem-poincare-duality-frobenius
\begin{proof}[{Proof}]\phantomsection\label{bp66}%
The pairing \(\beta_K\) is a composite of maps each of which commutes with pullback along Frobenius, so it is equivariant. Over a field the remaining statements then follow by linear algebra, because the pairing is perfect.

Cup product with mixed supports and evaluation are natural for pullback by Frobenius, and the isomorphism \(\gamma_{a_X}\) and the trace map are equivariant by Proposition \hyperref[bp62]{10.20}. Moreover, the contragredient Weil structure on \(K^\vee\) is chosen so that evaluation \(K\otimes K^\vee\to\Lambda_X\) is a morphism of Weil complexes, with the trivial structure on \(\Lambda_X\). It follows that the composite \(\beta_K\) is equivariant, which gives \(\langle Fx,Fy\rangle=\langle
x,y\rangle\) for the first pairing. The second pairing is obtained by moving the twist from the second factor to the values. Now geometric Frobenius acts on \(\Lambda(d)\) by \(q^{-d}\) and on \(\Lambda(-d)\) by \(q^d\), so for \(y\) in \(H^{2d-i}(X,K^\vee)\) the operator \(F\) on \(y\otimes\zeta\), with \(\zeta\) a generator of \(\Lambda(d)\), is \(q^{-d}(Fy)\otimes\zeta\). Removing the twist, we get the displayed identity \(\langle Fx,Fy\rangle=q^d\langle x,y\rangle\).

For field coefficients, Corollary \hyperref[bp40]{10.13} identifies \(H^{2d-i}(X,K^\vee)\), via the untwisted pairing, with the dual of \(H_c^i(X,K)\) tensored with \(\Lambda(-d)\). Moreover, Frobenius on \(H_c^i(X,K)\) is invertible, because \(\phi_K\) is an isomorphism and pullback along the automorphism \(\operatorname{Fr}_X\) is invertible on cohomology. In matrix terms, choose dual bases for the perfect pairing, let \(F_i\) be the matrix of \(F\) on \(H^i_c(X,K)\) and \(F_{2d-i}\) that on \(H^{2d-i}(X,K^\vee)\). Then the identity \(\langle Fx,Fy\rangle=q^d\langle x,y\rangle\) reads \({}^tF_iF_{2d-i}=q^d\). In other words, Frobenius on \(H^{2d-i}(X,K^\vee)\) is \(q^d\) times the inverse transpose of Frobenius on \(H_c^i(X,K)\). Since a matrix and its transpose have the same characteristic polynomial, this proves the eigenvalue statement with algebraic multiplicities. Finally, taking determinants of \(F_{2d-i}=q^d\,{}^t(F_i)^{-1}\) on spaces of dimension \(b\), we get \(\det F_{2d-i}=q^{db}(\det F_i)^{-1}\).
\end{proof}

% text: prose-10-24
\phantomsection\label{bp67}%
For \(X_0\) proper and geometrically connected and \(K=\Lambda_X\), Theorem \hyperref[bp46]{10.15} shows that \(F\) is the identity on \(H^0(X,\Lambda)=\Lambda\) and acts by \(q^d\) on \(H^{2d}(X,\Lambda)\simeq\Lambda(-d)\).

% object: example-duality-multiplicative-group (blueprint/chap-Verdier-and-Poincare-duality/sec-Frobenius-and-Poincare-duality/example-duality-multiplicative-group/_section.md)
\phantomsection\label{bp68}%
\begin{bpstatement}{Example~10.22}{ (Duality on the punctured line)}
Let \(X_0=\mathbf G_{m,\mathbf F_q}\), with coordinate \(t\), let \(X\) be its base change to \(k\), and let \(\Lambda=\overline{\mathbf Q}_\ell\), with the constant Weil sheaf. Since \(\operatorname{Pic}(X)=0\) and the units of \(k[t,t^{-1}]\) are the \(ct^m\), the Kummer sequence gives \(H^0(X,\Lambda)=\Lambda\) and \(H^1(X,\Lambda(1))=\Lambda\), generated by the Kummer class of \(t\) (Proposition \bpout{D.5}, with \(g=0\) and two punctures). Thus \(H^1(X,\Lambda)=\Lambda(-1)\), and the other groups vanish. For compact support, on the other hand, we write \(X=\mathbf P^1\setminus\{0,\infty\}\). The long exact sequence of \(j_!\Lambda_X\to\Lambda_{\mathbf P^1}\to\Lambda_{\{0,\infty\}}\to\) reads

\[
0\to H^0_c(X)\to\Lambda\xrightarrow{\ \text{diag}\ }\Lambda^2\to H^1_c(X)
\to H^1(\mathbf P^1)=0\to0\to H^2_c(X)\to H^2(\mathbf P^1)\to0,
\]

so \(H^0_c=0\), \(H^1_c=\Lambda^2/\Lambda\simeq\Lambda\) and \(H^2_c=H^2(\mathbf P^1,\Lambda)=\Lambda(-1)\). Poincaré duality with \(d=1\) pairs \(H^1_c\) with \(H^1\) and \(H^2_c\) with \(H^0\), and the dimensions match.

The Frobenius eigenvalues show the twist at work. On \(H^1_c\), which comes from the two rational points \(0\) and \(\infty\), Frobenius acts by \(1\), whereas on \(H^1=\Lambda(-1)\) it acts by \(q\). Their product is \(1\cdot q=q^d\), as Theorem \hyperref[bp65]{10.21} requires. Without the twist, duality would predict the eigenvalue \(1/1=1\) on \(H^1\) instead of \(q\). Likewise \(H^0\) and \(H^2_c\) have eigenvalues \(1\) and \(q\).
\end{bpstatement}

% object: example-running-example-duality (blueprint/chap-Verdier-and-Poincare-duality/sec-Frobenius-and-Poincare-duality/example-running-example-duality/_section.md)
\phantomsection\label{bp69}%
\begin{bpstatement}{Example~10.23}{ (The running example: duality)}
Let \(C_0\colon y^2=x^3-x\) over \(\mathbf F_3\) be the elliptic curve of the introduction, and \(C=C_0\times_{\mathbf F_3}k\). Since \(C\) has genus one, \(H^1(C,\mathbf Q_\ell)\) is a plane. Indeed, by the Kummer sequence \(H^1(C,\mu_{\ell^n})=J[\ell^n](k)\simeq(\mathbf Z/\ell^n)^{2g}\) with \(2g=2\) (Proposition \bpout{D.3} and Corollary \bpout{C.15}). Since \(k\) is algebraically closed, \(\mu_{\ell^n}\simeq\mathbf Z/\ell^n\), and passing to the limit and inverting \(\ell\) gives \(\dim H^1(C,\mathbf Q_\ell)=2\). For \(d=1\) the middle pairing on \(H^1\) is alternating and perfect, with values in \(\mathbf Q_\ell(-1)\) (Corollary \hyperref[bp52]{10.17}). By Theorem \hyperref[bp65]{10.21} it satisfies \(\langle Fx,Fy\rangle=3\langle x,y\rangle\). On a plane an alternating form is unique up to a scalar, so \(\langle Fx,Fy\rangle=\det(F)\langle
x,y\rangle\). Hence \(\det(F\mid H^1)=3\). If \(\alpha,\beta\) are the eigenvalues, then \(\alpha\beta=3\), and \(\det(1-TF\mid H^1)=1-(\alpha+\beta)T+3T^2\). This is the functional equation of a curve of genus one: replacing \(T\) by \(1/(3T)\) multiplies the polynomial by \((3T^2)^{-1}\).

Duality alone does not give \(|\alpha|=3^{1/2}\). For example, the pair \(\alpha=3\), \(\beta=1\), with \(\det(1-TF)=1-4T+3T^2\), consists of algebraic integers and satisfies \(\alpha\beta=3\). Here the count \(N_1=4\) gives \(\alpha+\beta=0\) and excludes it; in general only purity does.
\end{bpstatement}



\clearpage\phantomsection\hypertarget{bp-outside}{}%
\section*{Outside this export}
\addcontentsline{toc}{section}{Outside this export}
References marked \ensuremath{{}^\ast} are to parts of the blueprint this document leaves out:
\begin{description}
\item[Corollary 2.29] Continuity of cohomology, in Ⅰ The six operations (\bpslug{corollary-{}\allowbreak{}continuity-{}\allowbreak{}cohomology})
\item[Definition 2.45] Cohomological dimension, in Ⅰ The six operations (\bpslug{definition-{}\allowbreak{}cohomological-{}\allowbreak{}dimension})
\item[Section 2.13] Continuous cohomology, in Ⅰ The six operations (\bpslug{sec-{}\allowbreak{}Continuous-{}\allowbreak{}cohomology})
\item[Lemma 3.13] Flat models and generators, in Ⅰ The six operations (\bpslug{lemma-{}\allowbreak{}flat-{}\allowbreak{}models-{}\allowbreak{}generators})
\item[Proposition 5.15] Comparison with the étale derived category, in Ⅰ The six operations (\bpslug{proposition-{}\allowbreak{}comparison-{}\allowbreak{}etale-{}\allowbreak{}derived-{}\allowbreak{}category})
\item[Proposition 6.7] The adjunction for the projective line, in Ⅰ The six operations (\bpslug{proposition-{}\allowbreak{}adjunction-{}\allowbreak{}projective-{}\allowbreak{}line})
\item[Proposition 6.23] Projection formula with rational coefficients, in Ⅰ The six operations (\bpslug{proposition-{}\allowbreak{}projection-{}\allowbreak{}formula-{}\allowbreak{}rational-{}\allowbreak{}coefficients})
\item[Example 6.25] Local cohomology at the origin of the affine line, in Ⅰ The six operations (\bpslug{example-{}\allowbreak{}closed-{}\allowbreak{}support})
\item[Lemma 6.26] Projection formula with integral coefficients, in Ⅰ The six operations (\bpslug{lemma-{}\allowbreak{}projection-{}\allowbreak{}formula-{}\allowbreak{}integral-{}\allowbreak{}coefficients})
\item[Proposition 7.1] The six operations on $\Dctf$, in Ⅰ The six operations (\bpslug{proposition-{}\allowbreak{}six-{}\allowbreak{}operations-{}\allowbreak{}dctf})
\item[Proposition 7.10] The six operations and change of coefficients, in Ⅰ The six operations (\bpslug{proposition-{}\allowbreak{}six-{}\allowbreak{}operations-{}\allowbreak{}change-{}\allowbreak{}coefficients})
\item[Proposition 7.12] Isomorphisms reduce to finite coefficients, in Ⅰ The six operations (\bpslug{proposition-{}\allowbreak{}isomorphisms-{}\allowbreak{}reduce-{}\allowbreak{}finite-{}\allowbreak{}coefficients})
\item[Theorem 9.3] Smooth purity for all admissible coefficients, in Ⅱ Duality and the trace formula (\bpslug{theorem-{}\allowbreak{}smooth-{}\allowbreak{}purity-{}\allowbreak{}all-{}\allowbreak{}admissible-{}\allowbreak{}coefficients})
\item[Proposition 9.22] Compatibility of purity with base change, in Ⅱ Duality and the trace formula (\bpslug{proposition-{}\allowbreak{}compatibility-{}\allowbreak{}purity-{}\allowbreak{}base-{}\allowbreak{}change})
\item[Definition 9.24] Trace map, in Ⅱ Duality and the trace formula (\bpslug{definition-{}\allowbreak{}trace-{}\allowbreak{}map})
\item[Corollary 9.28] The dualizing complex of a smooth scheme, in Ⅱ Duality and the trace formula (\bpslug{corollary-{}\allowbreak{}dualizing-{}\allowbreak{}complex-{}\allowbreak{}smooth-{}\allowbreak{}scheme})
\item[Section 12.1] Frobenius correspondences, in Ⅱ Duality and the trace formula (\bpslug{sec-{}\allowbreak{}Frobenius-{}\allowbreak{}correspondences})
\item[Definition 12.2] Frobenius correspondences and their operators, in Ⅱ Duality and the trace formula (\bpslug{definition-{}\allowbreak{}frobenius-{}\allowbreak{}correspondences-{}\allowbreak{}their-{}\allowbreak{}operators})
\item[Lemma 12.3] Scheme Frobenius and geometric Frobenius, in Ⅱ Duality and the trace formula (\bpslug{lemma-{}\allowbreak{}scheme-{}\allowbreak{}frobenius-{}\allowbreak{}geometric-{}\allowbreak{}frobenius})
\item[Theorem 12.24] Grothendieck–Lefschetz trace formula, in Ⅱ Duality and the trace formula (\bpslug{theorem-{}\allowbreak{}grothendieck-{}\allowbreak{}lefschetz-{}\allowbreak{}trace-{}\allowbreak{}formula})
\item[Corollary 12.33] Functional equation, in Ⅱ Duality and the trace formula (\bpslug{corollary-{}\allowbreak{}functional-{}\allowbreak{}equation})
\item[Proposition 13.7] Weights under the usual operations, in Ⅲ Weights and the Riemann hypothesis (\bpslug{proposition-{}\allowbreak{}weights-{}\allowbreak{}under-{}\allowbreak{}usual-{}\allowbreak{}operations})
\item[Chapter 19] Proof of the Weil conjectures, in Ⅲ Weights and the Riemann hypothesis (\bpslug{chap-{}\allowbreak{}Proof-{}\allowbreak{}of-{}\allowbreak{}the-{}\allowbreak{}Weil-{}\allowbreak{}conjectures})
\item[Theorem A.2] Nagata's compactification theorem, in Appendices (\bpslug{theorem-{}\allowbreak{}nagata-{}\allowbreak{}s-{}\allowbreak{}compactification-{}\allowbreak{}theorem})
\item[Corollary C.15] Torsion of the Jacobian and field extension, in Appendices (\bpslug{corollary-{}\allowbreak{}torsion-{}\allowbreak{}jacobian-{}\allowbreak{}field-{}\allowbreak{}extension})
\item[Proposition D.3] Smooth proper curves, torsion prime to the characteristic, in Appendices (\bpslug{proposition-{}\allowbreak{}smooth-{}\allowbreak{}proper-{}\allowbreak{}curves-{}\allowbreak{}torsion-{}\allowbreak{}prime})
\item[Proposition D.5] Cohomology of a smooth affine curve, in Appendices (\bpslug{proposition-{}\allowbreak{}cohomology-{}\allowbreak{}smooth-{}\allowbreak{}affine-{}\allowbreak{}curve})
\end{description}
\clearpage\phantomsection%
\section*{References}
\addcontentsline{toc}{section}{References}
\bpbibitem{BBD82}{bp72}{A. Beilinson, J. Bernstein, and P. Deligne. \emph{Faisceaux pervers}, volume 100 of \emph{Astérisque}. 1982.}

\bpbibitem{BS15}{bp73}{B. Bhatt and P. Scholze. The pro-étale topology for schemes. \emph{Astérisque}, 369:99–201, 2015. \href{https://people.mpim-bonn.mpg.de/scholze/proetale.pdf}{Author's manuscript}.}

\bpbibitem{Con07}{bp74}{B. Conrad. Deligne's notes on Nagata compactifications. \emph{Journal of the Ramanujan Mathematical Society}, 22:205–257, 2007. \href{https://math.stanford.edu/\~conrad/papers/nagatafinal.pdf}{Author's manuscript}.}

\bpbibitem{Del74}{bp75}{P. Deligne. La conjecture de Weil. I. \emph{Publ. Math. IHÉS}, 43:273–307, 1974.}

\bpbibitem{Del80}{bp76}{P. Deligne. La conjecture de Weil. II. \emph{Publ. Math. IHÉS}, 52:137–252, 1980.}

\bpbibitem{Den22}{bp77}{C. Deninger. There is no “Weil-”cohomology theory with real coefficients for arithmetic curves. \emph{Ann. Sc. Norm. Super. Pisa Cl. Sci.}, 2022. To appear. arXiv:\href{https://arxiv.org/abs/2204.02714}{2204.02714}.}

\bpbibitem{Deu41}{bp78}{M. Deuring. Die Typen der Multiplikatorenringe elliptischer Funktionenkörper. \emph{Abh. Math. Sem. Univ. Hamburg}, 14:197–272, 1941.}

\bpbibitem{Eke90}{bp79}{T. Ekedahl. On the adic formalism. In \emph{The Grothendieck Festschrift, Vol. II}, volume 87 of \emph{Progress in Mathematics}, pages 197–218. Birkhäuser, 1990.}

\bpbibitem{Has36}{bp80}{H. Hasse. Zur Theorie der abstrakten elliptischen Funktionenkörper. I, II, III. \emph{J. Reine Angew. Math.}, 175:55–62, 69–88, 193–208, 1936.}

\bpbibitem{Hum72}{bp81}{J. E. Humphreys. \emph{Introduction to Lie Algebras and Representation Theory}, volume 9 of \emph{Graduate Texts in Mathematics}. Springer, 1972.}

\bpbibitem{Ill94}{bp82}{L. Illusie. Autour du théorème de monodromie locale. \emph{Astérisque}, 223:9–57, 1994. §4. \href{https://www.numdam.org/item/AST_1994__223__9_0/}{Original article}.}

\bpbibitem{Kat88}{bp83}{N. M. Katz. Travaux de Laumon. \emph{Astérisque}, 161–162:105–132, 1988. Séminaire Bourbaki, exp. 691.}

\bpbibitem{KW01}{bp84}{R. Kiehl and R. Weissauer. \emph{Weil conjectures, perverse sheaves and \(\ell\)-adic Fourier transform}, volume 42 of \emph{Ergebnisse der Mathematik}. Springer, 2001.}

\bpbibitem{Lau87}{bp85}{G. Laumon. Transformation de Fourier, constantes d'équations fonctionnelles et conjecture de Weil. \emph{Publ. Math. IHÉS}, 65:131–210, 1987.}

\bpbibitem{Mil86a}{bp86}{J. S. Milne. \href{https://www.jmilne.org/math/xnotes/AVs.pdf}{Abelian Varieties}. In \emph{Arithmetic Geometry}. 1986. Corrected online version.}

\bpbibitem{Mil86b}{bp87}{J. S. Milne. \href{https://www.jmilne.org/math/xnotes/JVs.pdf}{Jacobian Varieties}. In \emph{Arithmetic Geometry}. 1986. Corrected online version.}

\bpbibitem{Ran39}{bp88}{R. A. Rankin. Contributions to the theory of Ramanujan's function \(\tau(n)\) and similar arithmetical functions. II. \emph{Proc. Cambridge Philos. Soc.}, 35:357–372, 1939.}

\bpbibitem{SGA4}{bp70}{M. Artin, A. Grothendieck, and J.-L. Verdier, editors. \emph{Théorie des topos et cohomologie étale des schémas}, volumes 269, 270, 305 of \emph{Lecture Notes in Mathematics}. Springer, 1972–1973. SGA 4.}

\bpbibitem{SGA4½}{bp71}{P. Deligne. \emph{Cohomologie étale}, volume 569 of \emph{Lecture Notes in Mathematics}. Springer, 1977. SGA 4½; \emph{Théorèmes de finitude}, §3, and \emph{Rapport sur la formule des traces}, §3, and \emph{Sommes trigonométriques}, §1.}

\bpbibitem{SGA5}{bp89}{A. Grothendieck, editor. \emph{Cohomologie \(\ell\)-adique et fonctions \(L\)}, volume 589 of \emph{Lecture Notes in Mathematics}. Springer, 1977. SGA 5.}

\bpbibitem{SGA7II}{bp90}{P. Deligne and N. Katz, editors. \emph{Groupes de monodromie en géométrie algébrique}, volume 340 of \emph{Lecture Notes in Mathematics}. Springer, 1973. SGA 7 II; Exposé XIII, (2.1.7).}

\bpbibitem{Stacks}{bp91}{The Stacks Project Authors. \href{https://stacks.math.columbia.edu}{The Stacks Project}. Individual results are cited by their stable tags in the text.}

\bpbibitem{Var05}{bp92}{Y. Varshavsky. Lefschetz–Verdier trace formula and a generalization of a theorem of Fujiwara. 2005. Version 2. arXiv:\href{https://arxiv.org/abs/math/0505564}{math/0505564}. \href{https://arxiv.org/abs/math/0505564v2}{Author's preprint}.}

\bpbibitem{Wei48a}{bp93}{A. Weil. \emph{Sur les courbes algébriques et les variétés qui s'en déduisent}. Hermann, Paris, 1948.}

\bpbibitem{Wei48b}{bp94}{A. Weil. \emph{Variétés abéliennes et courbes algébriques}. Hermann, Paris, 1948.}

\bpbibitem{Wei49}{bp95}{A. Weil. Numbers of solutions of equations in finite fields. \emph{Bull. Amer. Math. Soc.}, 55:497–508, 1949.}

\end{document}
